\chapter{Method}\label{ch:method}
% =============================================================================

The method starts from \ac{DPCC} and changes two of its parts: the generative model and the projection.
The planning problem is formulated first, and the controller is then described from the plant upwards:
the two plants and the controllers that drive them, the generative models and their network, the
constraint projection, and the environments with their constraint sets.

\section{Problem Formulation}\label{sec:method:formalisation}

\subsection{The Control Problem}

The plant is a discrete-time system with state $\vect{s}_{t}\in\mathcal{S}$ and action
$\vect{a}_{t}\in\mathcal{A}$,
\begin{equation}\label{eq:method:plant}
  \vect{s}_{t+1} = f(\vect{s}_{t},\vect{a}_{t}) + \vect{w}_{t},
  \qquad \|\vect{w}_{t}\|_{2}\le\gamma,
\end{equation}
where $f$ is a known approximate model and $\vect{w}_{t}$ collects model mismatch. A goal
$g$ is reached when an indicator $\phi(\vect{s}_{t},g)$ equals one. Available is a dataset of
$N$ demonstrations produced by an unknown stochastic expert; the demonstrations are
\emph{multimodal}, meaning several distinct ways of reaching the goal are present, which is
why a generative model rather than a regressor is used. At test time the system must in
addition satisfy \emph{novel} constraints
\begin{equation}\label{eq:method:constraints}
  \vect{s}_{t}\in\mathcal{S}_{t}\subseteq\mathcal{S},
  \qquad \vect{a}_{t}\in\mathcal{A}_{t}\subseteq\mathcal{A},
  \qquad \forall t,
\end{equation}
which are not assumed to hold in the demonstrations. This is the problem statement of
\parencite[Sec.~3]{romer2025diffusion}, restated here because every claim in \autoref{ch:results} is made
against it.

\subsection{The Plan Variable}

At control step $t$ the model proposes a horizon of $H+1$ state--action pairs,
\begin{equation}\label{eq:method:plan}
  \vect{x}
  = \left(\vect{s}_{t:t+H|t},\,\vect{a}_{t:t+H|t}\right) \in \R^{n},
  \qquad n = (H+1)\,d,
\end{equation}
where $\vect{s}_{t+i|t}$ denotes the prediction for time $t+i$ made at time $t$, and $d$ is the
per-step transition width. The generative model is conditional on a context
$\vect{c} = (\vect{s}_{t},g)$, imposed by inpainting: the entries of $\vect{x}$ corresponding to
the current observation are overwritten with their known values after every sampler step, and
the same entries are zeroed in every training target so that no gradient is spent on them.
The feasible set is the intersection of the plan-level constraints with the dynamics rows,
\begin{equation}\label{eq:method:feasible}
  \mathcal{Z}_{f} = \left\{\vect{x} \;\middle|\;
    \vect{s}_{t'|t}\in\tilde{\mathcal{S}}_{t'},\;
    \vect{a}_{t'|t}\in\mathcal{A}_{t'},\;
    \vect{s}_{t'+1|t} = f(\vect{s}_{t'|t},\vect{a}_{t'|t}),\;
    \forall t'\right\},
\end{equation}
and the closed loop applies the first action $\vect{a}_{t|t}$ of the selected plan and re-solves
at $t+1$.

\subsection{Notation}

This field's papers overload $t$, $u$ and $x$: control time, transport time, the velocity field
and the control input all compete for the same letters. \autoref{tab:notation} fixes one
convention for the thesis. Two of its decisions carry through every chapter. The plan is
$\vect{x}$, not $\tau$ as in \parencite{romer2025diffusion}, which frees $\ftime$ for transport time;
and $\ftime$ runs \emph{data-at-one}, so the diffusion index $k$ of \autoref{sec:method:diffusion}
counts down to the data while the transport index $k$ of \autoref{sec:method:fm} counts up to it. Every
equation reproduced from a source that uses the opposite direction has been mapped, and the
mapping is noted at the point of use.

Three further collisions arise once the control stack is written down, and the table resolves them
by decoration rather than by renaming, so that each subfield keeps its own idiom. The generative
velocity field is $\vfield$ and the drone's linear velocity is $\dot{\vect{p}}$ --- never $\vect{v}$.
The body moment is $\vect{\tau}\sidx{b}$, always bold and always subscripted; an undecorated $\tau$
is always transport time. And $\matr{J}(\vect{q})$ is the manipulator Jacobian while
$\matr{J}\sidx{b}$ is the drone's inertia tensor; the two never appear in the same equation.

\begin{table}[htpb]
  \caption[Notation]{Notation used throughout. Vectors are bold lowercase, matrices bold
  uppercase, per the institute's convention.}
  \centering
  \label{tab:notation}
  \begin{tabularx}{\linewidth}{lLl}
    \toprule
    Symbol & Meaning & Range / type \\
    \midrule
    $t$ & control timestep (physical time) & $\mathbb{N}$ \\
    $H$ & planning horizon; a plan holds $H+1$ steps & $\mathbb{N}$ \\
    $\vect{s}_{t},\vect{a}_{t}$ & state, action & $\mathcal{S}$, $\mathcal{A}$ \\
    $\vect{s}_{t+i|t}$ & prediction for $t+i$ made at $t$ & \\
    $\vect{x}$ & the plan, flattened & $\R^{n}$ \\
    $\vect{c}$ & conditioning context $(\vect{s}_{t},g)$ & \\
    $\ftime$ & transport time; $0$ prior, $1$ data & $[0,1]$ \\
    $r$ & the second (start-side) transport time & $[0,\ftime]$ \\
    $h$ & interval length $\ftime-r$ & $[0,1]$ \\
    $k$ & sampler step index & $0,\dots,\nfe$ \\
    $\nfe$ & step budget $=$ network evaluations per plan: denoising steps of the reverse process
      for the diffusion baseline, fixed with its schedule at training; steps of the \ac{ODE} solver
      for the flow-based models, chosen at deployment & $\mathbb{N}$ \\
    $\beta_{k},\bar{\alpha}_{k},\sigma_{k}$ & diffusion schedule quantities & \\
    $\vfield_{\theta}$ & instantaneous velocity network & $\R^{n}$ \\
    $\afield_{\theta}$ & average velocity network & $\R^{n}$ \\
    $\vect{\epsilon}_{\theta}$ & noise-prediction network (baseline) & $\R^{n}$ \\
    $\mathcal{Z}_{f}$ & feasible set, dynamics included & $\subset\R^{n}$ \\
    $\proj{\mathcal{Z}_{f}}$ & projection \eqref{eq:bg:mpc:proj} & \\
    $\eta$ & constraint activation threshold & $[0,1]$ \\
    $B$ & number of candidate plans & $\mathbb{N}$ \\
    $\gamma$ & model-mismatch bound & $\R_{>0}$ \\
    \midrule
    \multicolumn{3}{l}{\emph{low-level control (\autoref{sec:method:deployment})}}\\
    $\vect{p}$, $\dot{\vect{p}}$ & measured position, linear velocity & $\R^{3}$ \\
    $\vect{p}\sidx{des}$ & commanded setpoint (the plan's integral) & $\R^{3}$ \\
    $\vect{q}$ & manipulator joint positions & $\R^{7}$ \\
    $\matr{J}(\vect{q})$ & manipulator Jacobian & $\R^{6\times 7}$ \\
    $\matr{R}$, $\vect{\omega}$ & body attitude, body rates & $SO(3)$, $\R^{3}$ \\
    $T$, $\vect{\tau}\sidx{b}$ & collective thrust, body moment & $\R$, $\R^{3}$ \\
    $\matr{J}\sidx{b}$ & body inertia (never the Jacobian) & $\R^{3\times3}$ \\
    $t\sidx{s}$ & sampling time of the dynamics model & $\R_{>0}$ \\
    \bottomrule
  \end{tabularx}
\end{table}

\section{System Overview}\label{sec:method:overview}

At every control step an observation conditions the generative model, the model proposes a plan, a
projection makes the plan admissible, and the first action is executed. The generative model and the
projection method are the two parts varied in this thesis. The generative models share the backbone
and its parameter count, the data, the normalisation, the horizon, the seeds, the projection and the
evaluation code; their training settings are given with the training protocol in
\autoref{sec:setup:protocol}. The projection methods share the feasible set, the solver, the activation
threshold and the number of candidate plans.
\hole{System-overview figure: observation $\rightarrow$ plan $\rightarrow$ projection
$\rightarrow$ action, with the two varied parts marked.}
\hole{Provenance table --- inherited, modified, new, at file granularity --- feeding
Appendix~\ref{app:repro}.} \srcnote{TARGET \S5.3}

% =============================================================================
%  PLANT DYNAMICS AND CONTROL
% =============================================================================
\section{Plant Dynamics and Control}\label{sec:method:dynamics}\label{sec:method:deployment}

Two plants are controlled: a manipulator, whose end-effector position the generative model plans, and
a quadrotor, whose position it plans. The projection uses a first-order model of the manipulator; the
quadrotor is described by its rigid-body dynamics. Underneath both sits a low-level controller, because
the generative model plans position increments and never an actuator command: the manipulator uses
differential \ac{IK} and a joint-space PD law, the quadrotor a cascaded geometric tracking controller
or, as a baseline, sampling-based predictive control.

\subsection{The Manipulator and Its First-Order Model}\label{sec:method:manipulator}

The projection of \autoref{sec:method:proj} needs a concrete dynamics model $f$ for its rows
\eqref{eq:bg:mpc:dyn}. The inherited system does not use an identified
model of the manipulator and its low-level controller; it uses a first-order (explicit Euler)
integrator over the commanded quantity, and pushes everything it omits into the mismatch term
$\vect{w}_{t}$ of \eqref{eq:method:plant}. For the two-dimensional manipulator benchmark, where the
state stacks a commanded and a measured planar position and the action is a planar increment, that
model is
\begin{equation}\label{eq:method:dpcc:euler}
  \vect{s}_{t+1}
  = \vect{s}_{t} + \begin{bmatrix}\vect{a}_{t}\trans & \vect{a}_{t}\trans\end{bmatrix}\trans t\sidx{s}
    + \vect{w}_{t},
\end{equation}
i.e.\ \emph{both} halves of the state are advanced by the \emph{same} action
\parencite[Sec.~6.1]{romer2025diffusion}. Three things follow, and all three are inherited by every
environment in this thesis.

\begin{itemize}
  \item \textbf{The model carries no knowledge of the robot.} $f$ knows nothing of the arm's
    kinematics, its impedance controller or its integration error. What it omits appears as the
    mismatch $\vect{w}_{t}$, whose bound $\gamma$ is estimated from unconstrained rollouts, and the
    tightening of \eqref{eq:method:dpcc:tighten} absorbs it. Feasibility is therefore a statement
    about the model and its mismatch bound, not about the plant.
  \item \textbf{$t\sidx{s} = 1$, because the action is a displacement and not a velocity.} The
    network's output is a per-step Cartesian increment, so the Euler update is already in the right
    units and no physical timestep enters the constraint. Setting $t\sidx{s}$ to the simulator's
    integration step --- a natural-looking mistake --- would rescale the entire feasible set.
    \srcnote{\texttt{aux\_repo/dpcc/config/projection\_eval.yaml:14--16}, whose comment reads
    \emph{``since $a = [\Delta x,\Delta y]$ and not $[v_{x},v_{y}]$''}; the same value and the same
    reasoning are carried in \texttt{config/uav\_projection.yaml:125--127} and
    \texttt{config/visual\_aligning\_eval.yaml}.}
  \item \textbf{Anchoring both the commanded and the measured channel is what couples the plan to
    the plant.} If only the measured position were constrained, the projector could return a plan
    whose commanded channel is unreachable; if only the commanded one were, the plan would be
    admissible in a space the robot never occupies. \autoref{sec:method:deployment} shows the same
    two-channel structure in all three environments, and \autoref{sec:method:proj} gives the row form.
\end{itemize}

\begin{remark}[A stricter alternative to the surrogate]\label{rem:eulermodel}
\eqref{eq:method:dpcc:euler} models the plant by an integrator: it carries no joint limits, no
impedance controller and no simulator integration error, and everything it omits is pushed into
$\gamma$ and absorbed by the tightening \eqref{eq:method:dpcc:tighten}. The alternative is to put the
\emph{real} kinematics inside the optimisation, and at least one contemporary system does exactly that.
\parencite{oelerich2026safeflowmpc} pairs a flow-matching planner with an online nonlinear program on
a seven-degree-of-freedom manipulator, and plans in \emph{joint space}: the generated trajectory is a
sequence of joint configurations, and the program carries the analytic forward kinematics, the
Jacobian and per-link collision geometry, solved at every flow-matching step. There is no surrogate
and no mismatch ball, because the constraint is imposed on the true kinematic chain.

The two designs sit at opposite ends of one trade. Exact kinematics inside the program give
feasibility on the real chain, and tie the program to one robot and the plan to that robot's joints.
The surrogate gives an end-effector plan and a projector that transfer unchanged from a planar arm to
a spatial one to a quadrotor, with feasibility that holds for the model and not for the plant. This
thesis takes the second end: one projection method, held fixed across two embodiments, is what RQ3 is about.
$\gamma$ is where the cost is booked.
\srcnote{Verified against the released code:
\texttt{aux\_repo/SafeFlowMPC/safe\_flow\_mpc/RobotModel/} carries the \texttt{iiwa.urdf} and
compiled CasADi functions for the forward kinematics, the Jacobian and ten per-link collision
positions; \texttt{train\_imitation\_learning.py} generates \texttt{n\_horizon} $\times$ 7 joint
trajectories. It plans \emph{joint configurations}, not motor torques.}
\end{remark}

\subsection{The Quadrotor}\label{sec:method:quadrotor}

The vehicle is the rigid body of \parencite[Eqs.~(2)--(5)]{lee2010geometric}: with mass $m$, body
inertia $\matr{J}\sidx{b}$, attitude $\matr{R}\in SO(3)$ and body rates $\vect{\omega}$,
\begin{equation}\label{eq:method:dep:plant}
  m\ddot{\vect{p}} = m g\,\unitvec{3} - T\,\matr{R}\unitvec{3},
  \qquad
  \dot{\matr{R}} = \matr{R}\hat{\vect{\omega}},
  \qquad
  \matr{J}\sidx{b}\dot{\vect{\omega}} + \vect{\omega}\times\matr{J}\sidx{b}\vect{\omega} = \vect{\tau}\sidx{b},
\end{equation}
where $\hat{(\cdot)}$ is the hat map.

The four rotor thrusts produce the collective thrust $T$ along the body axis $\matr{R}\unitvec{3}$ and
the body moment $\vect{\tau}\sidx{b}$. With four inputs for six degrees of freedom the vehicle is
underactuated: it translates only by tilting its thrust vector, and hover is an unstable equilibrium.
The generative model never commands these inputs; it plans positions, and the controller of
\autoref{sec:method:geometric} turns them into rotor thrusts.

\subsection{From Plan to Setpoint}\label{sec:method:setpoint}

In every case the generative model replans at each control step, emits a \emph{position increment},
and never emits an actuator command; everything below the increment is classical control --- a
differential \ac{IK} loop feeding a joint-space PD law on the two manipulator environments, and a
cascaded \emph{geometric} tracking controller, or optionally a sampling-based \ac{MJPC} tracker, on
the quadrotor.

Write $\Delta\vect{p}_{t}$ for the first action of the selected plan at control step $t$. All three
environments run the same two-line recursion,
\begin{align}
  \vect{p}_{\textrm{des},t+1} &= \vect{p}_{\textrm{des},t} + \Delta\vect{p}_{t},
    &&\vect{p}_{\textrm{des},0} = \vect{p}_{0},
    \label{eq:method:dep:setpoint}\\
  \vect{o}_{t} &= \left(\vect{p}_{\textrm{des},t},\;\vect{p}_{t},\;[\,\dot{\vect{p}}_{t}\,]\right),
    \label{eq:method:dep:obs}
\end{align}
where $\vect{o}_{t}$ is the observation the model conditions on at the next step and the bracketed
velocity block is present only for the aerial platform. Two properties follow.

\begin{itemize}
  \item \textbf{\eqref{eq:method:dep:setpoint} is a free-running integrator driven by the model's own
    past outputs.} $\vect{p}\sidx{des}$ is not read back from the plant. Plan error therefore
    accumulates rather than being corrected, which is the structural reason the aerial task is the
    fragile one and the reason a divergence guard is needed at all.
  \item \textbf{\eqref{eq:method:dep:obs} feeds the commanded setpoint back alongside the measured
    state}, so the model conditions on its own history. The gap
    $\vect{p}_{\textrm{des},t} - \vect{p}_{t}$ is exactly the tracking error the low-level controller is
    working on, and it is visible to the planner --- which is why the two channels are both
    constrained in \eqref{eq:method:proj:deriv}.
\end{itemize}

\begin{table}[htpb]
  \caption[Embodiment instantiation]{The same skeleton, instantiated three times. The plan width $d$
  is the width of the plan variable \eqref{eq:method:plan} and $\vect{o}$ its observation block;
  $t\sidx{s}$ is the Euler sampling time of \eqref{eq:method:dpcc:euler}; ``dynamics rows'' counts the
  channel pairs of \eqref{eq:method:proj:deriv}; the two intervals are simulated time
  \eqref{eq:method:dep:multirate}. The quantities measured on these loops are defined in
  \autoref{sec:setup:metrics}.}
  \label{tab:embodiments}
  \centering\small
  \begin{tabularx}{\linewidth}{lLLL}
    \toprule
     & D3IL-avoiding & D3IL-aligning & UAV-corridor, UAV-pillars, UAV-s-curve \\
    \midrule
    embodiment      & manipulator, planar   & manipulator, spatial      & quadrotor \\
    observation     & state, $4$-D          & 2 RGB views $+$ $6$-D state & state, $9$-D \\
    action          & $\Delta\vect{p}\in\R^{2}$ & $\Delta\vect{p}\in\R^{3}$ & $\Delta\vect{p}\in\R^{3}$ \\
    plan width $d$  & $6$                   & $9$                       & $12$ \\
    $\vect{o}$ layout & $(\vect{p}\sidx{des},\vect{p})$ & $(\vect{p}\sidx{des},\vect{p})$ & $(\vect{p}\sidx{des},\vect{p},\dot{\vect{p}})$ \\
    $t\sidx{s}$     & $1$                   & $1$                       & $1$ \\
    dynamics rows   & $4$                   & $6$                       & $6$ \\
    low level       & differential \ac{IK} $\rightarrow$ joint PD & same & geometric cascade, or \ac{MJPC} \\
    plan interval   & one plan per environment step & one plan per environment step & $0.03\,\textrm{s}$ simulated \\
    inner interval  & $n\sidx{sub}$ simulator substeps & $n\sidx{sub}$ simulator substeps & $0.01\,\textrm{s}$ simulated \\
    \bottomrule
  \end{tabularx}
\end{table}

\subsection{Manipulator: Differential Inverse Kinematics}\label{sec:method:ik}

The environment is commanded with an absolute Cartesian pose. The evaluation loop forms it from the
increment and the \emph{commanded} channel of the observation --- not from the measured one, which
is what makes \eqref{eq:method:dep:setpoint} the free-running recursion it is --- and holds the
end-effector orientation and, on the planar benchmark, the height fixed.
\srcnote{\texttt{aux\_repo/dpcc/scripts/eval.py:236--238}: \texttt{next\_pos\_des = action + obs[:2]}.}

Below that setpoint sits the differential \ac{IK} loop. Let $\vect{q}\in\R^{7}$ be the joint
positions of the Panda, $\matr{J}(\vect{q})$ the end-effector Jacobian, and let the task-space error be a
proportional term on position and orientation,
\begin{equation}\label{eq:method:dep:taskerr}
  \vect{\xi}(\vect{q})
  = \begin{bmatrix}
      k\sidx{p}\odot\operatorname{clip}\!\left(\vect{p}\sidx{des}-\vect{p}(\vect{q}),\,\pm\epsilon\sidx{p}\right)\\[2pt]
      k\sidx{q}\odot\operatorname{clip}\!\left(\vect{e}\sidx{rot}(\vect{q}),\,\pm\epsilon\sidx{r}\right)
    \end{bmatrix},
\end{equation}
with $\vect{p}(\vect{q})$ from forward kinematics and $\vect{e}\sidx{rot}$ the quaternion error to the
commanded orientation. The joint increment is the weighted, damped least-squares solution with a
null-space term that pulls the arm towards a rest posture,
\begin{align}
  \dot{\vect{q}}\sidx{null} &= k\sidx{null}\left(\vect{q}\sidx{rest}-\vect{q}\right),
    \label{eq:method:dep:null}\\
  \dot{\vect{q}}
    &= \matr{W}\matr{J}\trans
       \left(\matr{J}\matr{W}\matr{J}\trans + \lambda\Id{6}\right)^{-1}
       \left(\vect{\xi} - \matr{J}\dot{\vect{q}}\sidx{null}\right)
     + \dot{\vect{q}}\sidx{null},
    \label{eq:method:dep:dls}
\end{align}
after which the singular values of $\matr{J}\matr{W}\matr{J}\trans + \lambda\Id{6}$ are clipped into
$[\sigma\sidx{min},\sigma\sidx{max}]$ before the solve, $\|\dot{\vect{q}}\|$ is capped, and
$\vect{q} \leftarrow \operatorname{clip}(\vect{q}+\eta\sidx{ik}\dot{\vect{q}},\,\vect{q}\sidx{min},\vect{q}\sidx{max})$
is iterated $n\sidx{ik}$ times per control step. The resulting joint target, its finite-difference
velocity and acceleration are handed to a joint-space PD law, and \emph{that} is what produces the
torques the simulator integrates.

The parameters are inherited from the benchmark and are listed once because the loop is otherwise
not reproducible: $k\sidx{p} = (200,200,800)$, $k\sidx{q} = (30,30,30)$, $k\sidx{null} = 40$,
$\lambda = 10^{-12}$, $\matr{W} = \Id{7}$, $[\sigma\sidx{min},\sigma\sidx{max}] = [10^{-2},10^{2}]$,
$n\sidx{ik} = 3$, $\eta\sidx{ik} = 10^{-3}$, clips $\epsilon\sidx{p} = 0.01\,\textrm{m}$ and
$\epsilon\sidx{r} = 0.1$, and joint PD gains
$(120,120,120,120,50,30,10)$ / $(10,10,10,10,6,5,3)$.
\srcnote{\texttt{d3il/environments/d3il/d3il\_sim/controllers/IKControllers.py:134--323}
(\texttt{CartPosQuatImpedenceController}) and
\texttt{.../controllers/Config/mujoco\_controller\_config.gin:6--37}.}

Two consequences matter for the thesis. \eqref{eq:method:dep:dls} is a \emph{local} solve run three
times per step, so the arm does not exactly reach $\vect{p}\sidx{des}$ within one step --- the residual
is precisely the $\vect{w}_{t}$ that \eqref{eq:method:dpcc:euler} pushes into the mismatch bound and
\eqref{eq:method:dpcc:tighten} then absorbs. And the null-space term makes the map from setpoint to
joint configuration history-dependent, so two identical plans executed from different postures do not
produce identical joint trajectories, even though they produce the same Cartesian trace.

\subsection{Quadrotor: Cascaded Geometric Tracking Control}\label{sec:method:geometric}

The aerial chain is the same skeleton with two more layers underneath, and it runs on two intervals of
simulated time: a plan is queried every $n\sidx{dec}$-th physics step, the interval at which the expert
demonstrations were recorded, and the low-level controller runs at every physics step, with the
setpoint held constant in between,
\begin{equation}\label{eq:method:dep:multirate}
  n\sidx{dec}\,\Delta t\sidx{phys} = 0.03\,\textrm{s},
  \qquad
  \Delta t\sidx{phys} = 0.01\,\textrm{s}
  \;\Rightarrow\; n\sidx{dec} = 3 .
\end{equation}
The simulation waits for the planner, so these intervals hold in simulated time for any planning
time.
\srcnote{\texttt{mix\_uav\_test/eval\_mix\_uav.py:1610} (\texttt{decim}) and \texttt{:1728} (physics loop);
\texttt{uav\_expert\_data\_collect/dataset\_writer.py:31,73} (\texttt{DATASET\_HZ});
$\Delta t\sidx{phys}$ from \texttt{quadrotor\_modified.xml:4}.}

The controller is the \emph{position-controlled flight mode} of
\parencite[Eqs.~(19)--(23)]{lee2010geometric}, instantiated on \eqref{eq:method:dep:plant} in a
$z$-up frame with the thrust along $+\vect{b}_{3}$, and closed by the rotor allocation of
\parencite{mellinger2011minimum}. It runs in four steps.

\paragraph{Outer loop --- position to thrust vector.} With the tracking errors
$\vect{e}\sidx{p} = \vect{p}-\vect{p}\sidx{des}$ and
$\vect{e}\sidx{v} = \dot{\vect{p}}-\dot{\vect{p}}\sidx{des}$
\parencite[Eqs.~(17)--(18)]{lee2010geometric}, a diagonal PD with acceleration feed-forward gives
the required world-frame force
\begin{equation}\label{eq:method:dep:outer}
  \vect{a}\sidx{cmd} = -\matr{K}\sidx{p}\vect{e}\sidx{p} - \matr{K}\sidx{d}\vect{e}\sidx{v} + \ddot{\vect{p}}\sidx{des},
  \qquad
  \vect{F} = m\left(\vect{a}\sidx{cmd} + g\,\unitvec{3}\right),
\end{equation}
which is the bracketed term of \parencite[Eq.~(23)]{lee2010geometric}. The gain matrices are
diagonal, with a stiffer vertical channel.

\paragraph{Attitude extraction.} The rotor disk points along $\vect{F}$ and the remaining degree of
freedom is fixed by the commanded yaw $\psi\sidx{des}$, which pins the free axis left open by
\parencite[Eqs.~(22)--(23)]{lee2010geometric}:
\begin{equation}\label{eq:method:dep:att}
  \vect{b}_{3} = \frac{\vect{F}}{\|\vect{F}\|},\quad
  \vect{b}_{2} = \frac{\vect{b}_{3}\times\vect{x}\sidx{c}}{\|\vect{b}_{3}\times\vect{x}\sidx{c}\|},\quad
  \vect{b}_{1} = \vect{b}_{2}\times\vect{b}_{3},\quad
  \matr{R}\sidx{des} = \begin{bmatrix}\vect{b}_{1} & \vect{b}_{2} & \vect{b}_{3}\end{bmatrix},
\end{equation}
with $\vect{x}\sidx{c} = (\cos\psi\sidx{des},\sin\psi\sidx{des},0)\trans$.

\paragraph{Inner loop --- geometric attitude control.} The attitude error is the coordinate-free
error on $SO(3)$ of \parencite[Eq.~(21)]{lee2010geometric}, and the moment is its
\parencite[Eq.~(20)]{lee2010geometric}:
\begin{equation}\label{eq:method:dep:inner}
  \vect{e}\sidx{R} = \tfrac{1}{2}\left(\matr{R}\sidx{des}\trans\matr{R} - \matr{R}\trans\matr{R}\sidx{des}\right)^{\vee},
  \qquad
  \vect{\tau}\sidx{b} = -\matr{K}\sidx{R}\vect{e}\sidx{R} - \matr{K}_{\omega}\vect{e}_{\omega}
      + \vect{\omega}\times\matr{J}\sidx{b}\vect{\omega},
\end{equation}
where $(\cdot)^{\vee}$ is the inverse hat map. The commanded angular velocity is held at zero, so
$\vect{e}_{\omega} = \vect{\omega}$ and the attitude feed-forward of
\parencite[Eq.~(20)]{lee2010geometric} vanishes; with it goes the exponential-stability statement of
\parencite[Prop.~3]{lee2010geometric}, which assumes the full law.

\paragraph{Allocation.} The collective thrust is the world force projected on the current body-$z$
axis, floored so the vehicle cannot free-fall during recovery, and the wrench is distributed over the
four rotors by a fixed geometry matrix:
\begin{equation}\label{eq:method:dep:alloc}
  T = \max\!\left(\vect{F}\trans\vect{b}_{3},\;T\sidx{floor}\right),
  \qquad
  \matr{M}\vect{u} = \begin{bmatrix}T\\ \vect{\tau}\sidx{b}\end{bmatrix},
  \qquad
  \matr{M}_{:,i} = \begin{bmatrix}1 & r_{i,y} & -r_{i,x} & \kappa_{i}\end{bmatrix}\trans,
\end{equation}
with $\vect{r}_{i}$ the $i$-th rotor site position and $\kappa_{i}$ its yaw-torque gear ratio, both
read from the model file, so \eqref{eq:method:dep:alloc} holds for an arbitrary rotor placement
\parencite[Eq.~(1)]{lee2010geometric}\parencite{mellinger2011minimum}. When
$\vect{u} = \matr{M}^{-1}[T;\vect{\tau}\sidx{b}]$ leaves the actuator box, the command is rescaled
thrust-first: the mean is held and the torque component is shrunk by the largest factor that keeps
every rotor in range, but never below one half. Parameters of record:
$\matr{K}\sidx{p} = \operatorname{diag}(4,4,8)$, $\matr{K}\sidx{d} = \operatorname{diag}(3,3,4)$,
$\matr{K}\sidx{R} = \operatorname{diag}(70,70,4)$, $\matr{K}_{\omega} = \operatorname{diag}(2.5,2.5,1)$,
$u\sidx{hover} = mg/4$, $u\sidx{max} = \max(2u\sidx{hover},6)$, $u\sidx{min} = 0$,
$T\sidx{floor} = 0.1\,mg$, $g = 9.81\,\textrm{m}/\textrm{s}^{2}$; $m$ and $\matr{J}\sidx{b}$ are
read from the vehicle model.
The loop has no integral term: the four gains act on position, velocity, attitude and angular rate,
and the inner loop is the coordinate-free error on $SO(3)$ of \eqref{eq:method:dep:inner}, closed by
the static allocation \eqref{eq:method:dep:alloc}. The thesis therefore calls it a cascaded geometric
tracking controller; \texttt{pid} in the run tags is an artefact token.
\srcnote{\texttt{uav\_env\_test/flight\_controller.py:33--155} (class \texttt{CascadedPID}); the name
mapping is in Auxiliary/Naming/TRANSLATION\_20260914\_dev\_jargon\_to\_scientific.md \S7.}

\paragraph{The velocity setpoint is synthesised, not planned.} \eqref{eq:method:dep:outer} needs
$\dot{\vect{p}}\sidx{des}$, which the plan does not supply --- it supplies displacements. Three
policies are implemented and the choice is a recorded experimental factor, not an implementation
detail:
\begin{equation}\label{eq:method:dep:vdes}
  \dot{\vect{p}}\sidx{des} =
  \begin{cases}
    \dfrac{\Delta\vect{p}_{t}}{n\sidx{dec}\,\Delta t\sidx{phys}} & \text{default --- the increment read as a mean velocity,}\\
    \vect{0} & \text{stop-and-go --- brake to rest at every plan step,}\\
    \dfrac{\Delta\vect{p}_{t}}{\|\Delta\vect{p}_{t}\|}\,\bar{v} & \text{constant-speed, with }
      \bar{v} = \overline{\|\Delta\vect{p}\|}\,/\,(n\sidx{dec}\,\Delta t\sidx{phys})\text{ from the dataset.}
  \end{cases}
\end{equation}
The third is timing-free by construction and calibrates itself to whatever data it is given; its
magnitude is the mean of the first policy, so the three differ in \emph{profile}, not in average
speed. The first reads the increment as a mean velocity over an assumed constant interval between
plan queries, so the commanded speed moves with the arrival time of the plan.

Every quadrotor result reported in this thesis is flown with the second policy, braking to rest at
every plan step; the other two are implemented and were measured, and are not the configuration of
record. With $\dot{\vect{p}}\sidx{des} = \vect{0}$ the outer loop of \eqref{eq:method:dep:outer} moves
the vehicle on its position error alone, so the setpoint of \eqref{eq:method:dep:setpoint} leads the
vehicle by a standing lag that grows with the commanded speed. \autoref{sec:res:uav} reports what
that lag costs.
\srcnote{\texttt{mix\_uav\_test/eval\_mix\_uav.py:1551--1563, 1822--1832}; policy of record and its
measured lag from Gen11/Epoch8\_UAV\_Mjpc\_thrust\_control/U3\_v\_des\_Patch/PLAN\_pid\_const\_v.md and
Gen15/Study/STUDY\_20260901\_mf\_unguided\_failure\_uav\_pillars.md \S5.}

\subsection{Sampling-Based Predictive Control}\label{sec:method:mjpc}

A sampling-based predictive controller is
implemented behind the same interface, so it can be swapped for the geometric cascade without
touching the plan layer. It is the \emph{Predictive Sampling} planner of \parencite{howell2022predictive}, used
as published; only the task is ours. Three pieces define it.

\emph{The objective.} The planner minimises the finite-horizon sum
$J = \sum_{i} c(\vect{x}_{i},\vect{u}_{i})$ over a rollout of the true simulator dynamics
\parencite[Eq.~(3)]{howell2022predictive}, with the running cost a weighted sum of twice-differentiable
norms of \emph{residuals} --- quantities that are small when the task is solved
\parencite[Eq.~(4)]{howell2022predictive}. The task defined here has two residual terms, both under a
quadratic norm and at risk-neutral setting ($R = 0$ in
\parencite[Eq.~(5)]{howell2022predictive}, so the risk transform is the identity):
\begin{equation}\label{eq:method:dep:mjpc}
  J = \sum_{i=1}^{H\sidx{mjpc}}
      \Big(\underbrace{\left\|\vect{p}_{i}-\vect{p}\sidx{des}\right\|_{2}^{2}}_{r_{1},\;w_{1}=1}
      + \underbrace{w_{v}\left\|\dot{\vect{p}}_{i}\right\|_{2}^{2}}_{r_{2},\;w_{2}=w_{v}}\Big),
\end{equation}
with $\vect{p}\sidx{des}$ passed in as the instruction. The velocity residual is not decoration: it
plays the role of the geometric cascade's derivative gain, and without it the planner reaches the setpoint at
full thrust and overshoots --- minimum cumulative position error being achieved by arriving as fast
as possible.

\emph{The plan and its parameterisation.} The plan is a time-indexed spline over control knots
$\vect{\theta}$ rather than a raw action sequence \parencite[\S3.3]{howell2022predictive}; the
zero-order-hold interpolant is used here, with one knot per horizon step, so the spline degenerates
to the action sequence itself and the parameterisation buys nothing but interface uniformity.

\emph{The search.} Each call performs a fixed number of iterations of Algorithm~4 of
\parencite{howell2022predictive}: draw $N-1$ perturbations
$\vect{\theta}_{i} \sim \Normal(\vect{\theta},\sigma^{2})$, roll out all $N$ candidates --- the
incumbent included --- from the measured state, and keep
$\vect{\theta} \leftarrow \arg\min_{i} J(\vect{\theta}^{(i)})$. It is derivative-free and zero-order;
its authors present it as an elementary baseline. Between control steps the plan is shifted forward
and re-sampled, which is the warm-start of
\parencite[Alg.~1]{howell2022predictive} and the reason a few iterations per step suffice.

\emph{One deliberate departure.} Contacts are disabled in the planner's copy of the model. The
tracker therefore plans against the vehicle's free-flight dynamics only, and obstacle avoidance
remains entirely the projector's job. This is not a shortcut: it keeps the constraint mechanism a
single identifiable component, so that swapping the tracker cannot silently change how constraints
are enforced, so the swap stays a controlled comparison.
Defaults: $N = 16$, $H\sidx{mjpc} = 0.3\,\textrm{s}$ of physics steps, five improvement iterations
per call, $\sigma = 0.3$, $w_{v} = 0.1$.
\srcnote{\texttt{mix\_uav\_test/mjpc\_tracker.py:70--155}, which wraps
\texttt{mujoco\_mpc.mjx.predictive\_sampling} unmodified.}

\hole{State once, in \autoref{sec:setup:protocol}, which controller produced each reported UAV
number. The apparatus notes record that every dataset and every reported result to date uses the
geometric cascade, and that the predictive controller was built and studied but is not the source of a
reported number --- confirm this against the run ledger before the claim is made in print.}

% =============================================================================
%  GENERATIVE MODELS
% =============================================================================
\section{Generative Models}\label{sec:method:engine}

The diffusion model of \ac{DPCC} is the baseline. Instantaneous-velocity matching, analytic
average-velocity matching and consistency-interpolated average-velocity matching replace it
behind one interface, each with its own training target and integration loop.

\subsection{Diffusion and the DPCC Baseline}\label{sec:method:diffusion}

Let $\vect{x} \in \R^{n}$ be the flattened plan of \eqref{eq:method:plan}. A denoising diffusion probabilistic model
\parencite{sohl2015deep,ho2020denoising} introduces latent variables
$\vect{x}^{1},\dots,\vect{x}^{\nfe}$ through a fixed Gaussian forward process
\begin{equation}\label{eq:bg:ddpm:forward}
  q\!\left(\vect{x}^{k} \mid \vect{x}^{k-1}\right)
  = \Normal\!\left(\sqrt{1-\beta_{k}}\,\vect{x}^{k-1},\; \beta_{k}\Id{n}\right),
  \qquad k = 1,\dots,\nfe ,
\end{equation}
whose step sizes $\beta_{1:\nfe} \in (0,1)^{\nfe}$ are set by a \emph{noise schedule}. Because
the transitions are Gaussian, the marginal is available in closed form,
\begin{equation}\label{eq:bg:ddpm:marginal}
  q\!\left(\vect{x}^{k} \mid \vect{x}^{0}\right)
  = \Normal\!\left(\sqrt{\bar{\alpha}_{k}}\,\vect{x}^{0},\;(1-\bar{\alpha}_{k})\Id{n}\right),
  \qquad
  \alpha_{k} = 1-\beta_{k}, \quad \bar{\alpha}_{k} = \prod_{i=1}^{k}\alpha_{i},
\end{equation}
and the schedule is chosen so that $q(\vect{x}^{\nfe}\mid\vect{x}^{0}) \approx \Normal(\vect{0},\Id{n})$.
Generation reverses \eqref{eq:bg:ddpm:forward} with a learned Markov chain
\begin{equation}\label{eq:bg:ddpm:reverse}
  p_{\theta}\!\left(\vect{x}^{k-1}\mid\vect{x}^{k},k,\vect{c}\right)
  = \Normal\!\left(\vect{\mu}_{\theta}\!\left(\vect{x}^{k},k,\vect{c}\right),\;\sigma_{k}^{2}\Id{n}\right),
  \qquad p(\vect{x}^{\nfe}) = \Normal(\vect{0},\Id{n}),
\end{equation}
with $\vect{c}$ the conditioning context and the variance pinned to the forward posterior,
$\sigma_{k}^{2} = \beta_{k}(1-\bar{\alpha}_{k-1})/(1-\bar{\alpha}_{k})$. The mean is learned
indirectly by predicting the noise that was added \parencite{ho2020denoising}, in the form
restated by \parencite[Sec.~4]{romer2025diffusion}:
\begin{equation}\label{eq:bg:ddpm:loss}
  \mathcal{L}(\theta)
  = \E_{k \sim \Unif\{1,\nfe\},\; \vect{x}^{0}\sim q,\; \vect{\epsilon}\sim\Normal(\vect{0},\Id{n})}
    \left\|\vect{\epsilon} - \vect{\epsilon}_{\theta}\!\left(
      \sqrt{\bar{\alpha}_{k}}\,\vect{x}^{0} + \sqrt{1-\bar{\alpha}_{k}}\,\vect{\epsilon},\;k\right)\right\|_{2}^{2}.
\end{equation}
At sampling time the noise prediction is turned back into a mean in two steps --- an estimate of
the clean sample, and the forward posterior evaluated at that estimate:
\begin{align}
  \hat{\vect{x}}^{0}\!\left(\vect{x}^{k},k\right)
    &= \tfrac{1}{\sqrt{\bar{\alpha}_{k}}}\,\vect{x}^{k}
       - \sqrt{\tfrac{1}{\bar{\alpha}_{k}}-1}\;\vect{\epsilon}_{\theta}\!\left(\vect{x}^{k},k,\vect{c}\right),
       \label{eq:bg:ddpm:x0} \\[2pt]
  \vect{\mu}_{\theta}\!\left(\vect{x}^{k},k,\vect{c}\right)
    &= \frac{\beta_{k}\sqrt{\bar{\alpha}_{k-1}}}{1-\bar{\alpha}_{k}}\,\hat{\vect{x}}^{0}
     + \frac{(1-\bar{\alpha}_{k-1})\sqrt{\alpha_{k}}}{1-\bar{\alpha}_{k}}\,\vect{x}^{k}.
       \label{eq:bg:ddpm:mean}
\end{align}
\srcnote{\eqref{eq:bg:ddpm:x0}--\eqref{eq:bg:ddpm:mean} are the two functions
\texttt{predict\_start\_from\_noise} and \texttt{q\_posterior} in
\texttt{aux\_repo/dpcc/diffuser/models/diffusion.py:97--117}; they are what the baseline
actually executes, and \autoref{sec:method:engine} replaces exactly this pair.}

Three properties are carried forward. First, \emph{$\nfe$ is a training-time constant}: it
indexes the schedule $\beta_{1:\nfe}$ that \eqref{eq:bg:ddpm:forward} and
\eqref{eq:bg:ddpm:marginal} are defined by, so a model trained at $\nfe = 20$ has no
well-defined $2$-step sampler. Reducing $\nfe$ after training requires a different sampler
--- an implicit or distilled one \parencite{song2021denoising} --- and therefore a different
model, not a different argument. The configuration files show the same asymmetry: the baseline's
step count is a \emph{training} key, the flow-matching models' an \emph{inference} key. \srcnote{\texttt{config/}:
\texttt{n\_diffusion\_steps} sits in the training block, \texttt{flow\_steps\_v3} in the plan
block.}
Second, the reverse chain is \emph{stochastic}: each step injects fresh noise of scale
$\sigma_{k}$, so two plans drawn from the same context differ. Third, the whole construction
is a chain of Gaussians, which is what allows a constraint mechanism to be inserted between
its steps (\autoref{sec:method:dpcc}).

The instance this thesis is measured against uses the cosine schedule of
\parencite{nichol2021improved} with $\nfe = 20$ and a temporal U-Net denoiser
\parencite{janner2022planning}, conditioned on the current state by inpainting rather than by
classifier or classifier-free guidance \parencite{dhariwal2021diffusion,ho2022classifier}, following
\parencite[Sec.~6.1]{romer2025diffusion}.

\subsection{Flow Matching}\label{sec:method:fm}

\subsubsection*{Transport, and the field that generates it}

Let $\ftime \in [0,1]$ be the \emph{transport time}, with $\ftime = 0$ at the prior and
$\ftime = 1$ at the data --- the convention used throughout this thesis and by this work's
implementation.\footnote{The opposite convention (noise at $\ftime=1$) is equally common;
\parencite{geng2025meanflow} uses it. Every equation reproduced here has been mapped into the
data-at-one convention, and the sign changes this induces are noted where they occur.}
A time-dependent velocity field $\vfield_{\ftime}:\R^{n}\to\R^{n}$ defines a flow
$\Phi$ through the ordinary differential equation
\begin{equation}\label{eq:bg:fm:ode}
  \frac{\dd}{\dd\ftime}\,\vect{x}_{\ftime} = \vfield_{\ftime}(\vect{x}_{\ftime}),
  \qquad \vect{x}_{0} \sim p_{0},
\end{equation}
and the pushforward $p_{\ftime} = (\Phi_{0\to\ftime})_{\sharp}p_{0}$ is the probability path
it generates, characterised by the continuity equation
$\partial_{\ftime}p_{\ftime} + \nabla\!\cdot(p_{\ftime}\vfield_{\ftime}) = 0$
\parencite[Sec.~2]{lipman2023flow}. Sampling is therefore integration: draw $\vect{x}_{0}\sim p_{0}$ and
solve \eqref{eq:bg:fm:ode} to $\ftime = 1$.

\subsubsection*{The objective, and why it is tractable}

Flow matching \parencite[Eq.~(5)]{lipman2023flow} regresses a network $\vfield_{\ftime}^{\theta}$ onto the
field that generates a prescribed path:
\begin{equation}\label{eq:bg:fm:fm}
  \mathcal{L}\sidx{FM}(\theta)
  = \E_{\ftime\sim\Unif[0,1],\;\vect{x}\sim p_{\ftime}}
    \left\|\vfield_{\ftime}^{\theta}(\vect{x}) - \vfield_{\ftime}(\vect{x})\right\|_{2}^{2}.
\end{equation}
This is not computable: the marginal field $\vfield_{\ftime}$ is an integral over the data
distribution. The construction that makes it usable is to build the path as a mixture of
\emph{per-sample} paths. Conditioning on a data point $\vect{x}_{1}$, one designs
$p_{\ftime}(\vect{x}\mid\vect{x}_{1})$ and a conditional field
$\vfield_{\ftime}(\vect{x}\mid\vect{x}_{1})$ that generates it; the marginals are then
\begin{equation}\label{eq:bg:fm:marginal}
  p_{\ftime}(\vect{x}) = \int p_{\ftime}(\vect{x}\mid\vect{x}_{1})\,q(\vect{x}_{1})\,\dd\vect{x}_{1},
  \qquad
  \vfield_{\ftime}(\vect{x}) = \int \vfield_{\ftime}(\vect{x}\mid\vect{x}_{1})
    \frac{p_{\ftime}(\vect{x}\mid\vect{x}_{1})\,q(\vect{x}_{1})}{p_{\ftime}(\vect{x})}\,\dd\vect{x}_{1},
\end{equation}
and the marginal field so defined does generate the marginal path
\parencite[Theorem~1]{lipman2023flow}. The \emph{conditional} objective
\begin{equation}\label{eq:bg:fm:cfm}
  \mathcal{L}\sidx{CFM}(\theta)
  = \E_{\ftime\sim\Unif[0,1],\;\vect{x}_{1}\sim q,\;\vect{x}\sim p_{\ftime}(\cdot\mid\vect{x}_{1})}
    \left\|\vfield_{\ftime}^{\theta}(\vect{x}) - \vfield_{\ftime}(\vect{x}\mid\vect{x}_{1})\right\|_{2}^{2}
\end{equation}
is an unbiased surrogate: it has the same gradient in $\theta$ as \eqref{eq:bg:fm:fm}
\parencite[Theorem~2]{lipman2023flow}. Training is therefore \emph{simulation-free} --- no
trajectory is ever integrated to form a target.

\subsubsection*{The linear path}

This thesis uses the linear (optimal-transport displacement, ``rectified'')
path \parencite{lipman2023flow,liu2023flow}, with the noise--data coupling drawn independently:
\begin{equation}\label{eq:bg:fm:path}
  \vect{x}_{\ftime} = (1-\ftime)\,\vect{x}_{0} + \ftime\,\vect{x}_{1},
  \qquad
  \vfield_{\ftime}(\vect{x}_{\ftime}\mid \vect{x}_{0},\vect{x}_{1}) = \vect{x}_{1} - \vect{x}_{0},
\end{equation}
i.e.\ the special case $\sigma\sidx{min} = 0$ of \parencite[Eq.~(20)--(23)]{lipman2023flow}, read in
the data-at-one direction. Substituting \eqref{eq:bg:fm:path} into \eqref{eq:bg:fm:cfm} gives
the objective this work trains:
\begin{equation}\label{eq:bg:fm:loss}
  \mathcal{L}\sidx{FM}(\theta)
  = \E_{\ftime,\;\vect{x}_{1}\sim q,\;\vect{x}_{0}\sim p_{0}}
    \left\|\vfield_{\theta}\!\left((1-\ftime)\vect{x}_{0} + \ftime\vect{x}_{1},\;\ftime,\;\vect{c}\right)
      - \left(\vect{x}_{1}-\vect{x}_{0}\right)\right\|_{2}^{2}.
\end{equation}
The regression target is a difference of two draws. Nothing in it depends on a schedule, and
nothing in it depends on how many steps the sampler will later take.

\subsubsection*{Sampling, and the two consequences that matter here}

With $\vfield_{\theta}$ trained, a plan is produced by integrating \eqref{eq:bg:fm:ode}. With
the explicit Euler scheme at a uniform grid $\ftime_{k} = k/\nfe$,
\begin{equation}\label{eq:bg:fm:euler}
  \vect{x}_{\ftime_{k+1}} = \vect{x}_{\ftime_{k}}
    + \tfrac{1}{\nfe}\,\vfield_{\theta}\!\left(\vect{x}_{\ftime_{k}},\ftime_{k},\vect{c}\right),
  \qquad k = 0,\dots,\nfe-1,
\end{equation}
which costs exactly $\nfe$ network evaluations. This count, the \ac{NFE}, is the step budget used
throughout the thesis. Two consequences follow.

\emph{The budget moved.} $\nfe$ appears in \eqref{eq:bg:fm:euler} and nowhere in
\eqref{eq:bg:fm:loss}. It is a solver argument, chosen after training and changeable per
deployment, in contrast to \eqref{eq:bg:ddpm:forward}--\eqref{eq:bg:ddpm:loss} where it
indexes the schedule the model was fitted to. What limits small $\nfe$ is now discretisation
error of \eqref{eq:bg:fm:ode}, which is a property of the field's curvature rather than of the
training objective --- and that is the quantity analytic average-velocity matching addresses
(\autoref{sec:method:meanflow}).

\emph{The sampler became deterministic.} \eqref{eq:bg:fm:euler} injects no noise: the plan is
a function of $\vect{x}_{0}$ and $\vect{c}$ alone. This removes a source of run-to-run variance
from a controller that replans at every timestep, and it changes what ``the same plan'' means
when comparing projection methods: two runs started from the same $\vect{x}_{0}$ differ only in the
projection, so the comparison in \autoref{ch:results} is controlled. The candidate plans of
\autoref{sec:method:selection} therefore differ in $\vect{x}_{0}$, not in sampler noise.

\subsubsection*{Implementation}

Instantaneous-velocity matching trains on exactly this objective. Training draws a data plan $\vect{x}_{1}$, a prior sample $\vect{x}_{0}$ and a transport time
$\ftime$; forms the interpolant and regresses the velocity, exactly
\eqref{eq:bg:fm:path}--\eqref{eq:bg:fm:loss}:
\begin{equation}\label{eq:method:engine:fm}
  \vect{x}_{\ftime} = (1-\ftime)\vect{x}_{0} + \ftime\vect{x}_{1},
  \qquad
  \mathcal{L}(\theta) = \E\left\|
     \vfield_{\theta}(\vect{x}_{\ftime},\ftime,\vect{c}) - (\vect{x}_{1}-\vect{x}_{0})
  \right\|_{2}^{2}.
\end{equation}
Sampling is \eqref{eq:bg:fm:euler}. Relative to the baseline this deletes the schedule buffers,
the noise-to-sample conversion \eqref{eq:bg:ddpm:x0} and the posterior mean
\eqref{eq:bg:ddpm:mean}, and it deletes the noise injected at every reverse step. The
conditioning mask, the loss reduction and the backbone are untouched.
\srcnote{\texttt{flow\_matcher\_v3/models/diffusion.py:273--299}
(\texttt{q\_sample}, \texttt{p\_losses}).}

\subsection{Analytic Average-Velocity Matching}\label{sec:method:meanflow}

\subsubsection*{Average velocity}

The few-step family changes \emph{what the network predicts}. Following
\parencite[Eq.~(3)]{geng2025meanflow}, define the \emph{average velocity} over an interval
$[r,\ftime] \subseteq [0,1]$ as the displacement divided by the interval length,
\begin{equation}\label{eq:bg:mf:def}
  \afield(\vect{x}_{\ftime},r,\ftime)
  \;\triangleq\; \frac{1}{\ftime-r}\int_{r}^{\ftime}\vfield(\vect{x}_{\varsigma},\varsigma)\,\dd\varsigma ,
\end{equation}
in contrast to the instantaneous velocity $\vfield$ of \autoref{sec:method:fm}. The field $\afield$ is
induced by $\vfield$; it involves no network. It satisfies the boundary condition
$\lim_{r\to\ftime}\afield = \vfield$, and, by additivity of the integral, the consistency
relation
$(\ftime-r)\afield(\vect{x}_{\ftime},r,\ftime) = (s-r)\afield(\vect{x}_{s},r,s)
+ (\ftime-s)\afield(\vect{x}_{\ftime},s,\ftime)$ for any intermediate $s$
\parencite[Sec.~4.1]{geng2025meanflow}. The point of \eqref{eq:bg:mf:def} is immediate: if a network
approximates $\afield$ well, then
\begin{equation}\label{eq:bg:mf:onestep}
  \vect{x}_{1} = \vect{x}_{0} + \afield_{\theta}(\vect{x}_{0},0,1)
\end{equation}
traverses the entire path in a single evaluation. No integral is approximated at inference;
the interval is a network input.

\subsubsection*{The average-velocity identity}

\eqref{eq:bg:mf:def} cannot be used as a regression target directly, because forming it
requires integrating during training. The construction of \parencite[Eqs.~(4)--(6)]{geng2025meanflow}
removes the
integral. Write \eqref{eq:bg:mf:def} as
$(\ftime-r)\,\afield(\vect{x}_{\ftime},r,\ftime) = \int_{r}^{\ftime}\vfield\,\dd\varsigma$ and
differentiate both sides with respect to $\ftime$ at fixed $r$; the product rule on the left
and the fundamental theorem of calculus on the right give the \emph{average-velocity identity}
\begin{equation}\label{eq:bg:mf:identity}
  \underbrace{\afield(\vect{x}_{\ftime},r,\ftime)}_{\text{average velocity}}
  \;=\;
  \underbrace{\vfield(\vect{x}_{\ftime},\ftime)}_{\text{instantaneous velocity}}
  \;-\;(\ftime-r)\,
  \underbrace{\frac{\dd}{\dd\ftime}\afield(\vect{x}_{\ftime},r,\ftime)}_{\text{total time derivative}} .
\end{equation}
The identity is equivalent to the definition \eqref{eq:bg:mf:def} and contains no neural
network. The total derivative expands, using $\dd\vect{x}_{\ftime}/\dd\ftime = \vfield$ and
$\dd r/\dd\ftime = 0$, into
\begin{equation}\label{eq:bg:mf:jvp}
  \frac{\dd}{\dd\ftime}\afield(\vect{x}_{\ftime},r,\ftime)
  = \vfield(\vect{x}_{\ftime},\ftime)\,\partial_{\vect{x}}\afield + \partial_{\ftime}\afield ,
\end{equation}
which is exactly a \ac{JVP} of $\afield$ with the tangent
$(\vfield,\,0,\,1)$ --- computable in one forward-mode pass, at a cost comparable to one
backward pass \parencite[Eqs.~(7)--(8)]{geng2025meanflow}.

Parameterising $\afield$ by a network $\afield_{\theta}$ and substituting the conditional
velocity of \eqref{eq:bg:fm:path} for the marginal one, as in \autoref{sec:method:fm}, yields the
objective
\begin{equation}\label{eq:bg:mf:loss}
  \mathcal{L}\sidx{avg}(\theta)
  = \E\left\|\afield_{\theta}(\vect{x}_{\ftime},r,\ftime)
      - \sg\!\left(\afield\sidx{tgt}\right)\right\|_{2}^{2},
  \qquad
  \afield\sidx{tgt} = \vfield_{\ftime} - (\ftime-r)\left(
     \vfield_{\ftime}\,\partial_{\vect{x}}\afield_{\theta} + \partial_{\ftime}\afield_{\theta}\right),
\end{equation}
where $\sg$ is the stop-gradient. The stop-gradient is what keeps the method first-order: the
\ac{JVP} enters only as a constant target, so no differentiation through it is
required \parencite[Eqs.~(9)--(11)]{geng2025meanflow}.

Three properties of \eqref{eq:bg:mf:loss} are used later in this thesis.

\begin{itemize}
  \item \textbf{It contains instantaneous-velocity matching as the case $r = \ftime$.} The second term of
    $\afield\sidx{tgt}$ carries the factor $(\ftime-r)$, so at $r=\ftime$ the target collapses to
    $\vfield_{\ftime}$ and \eqref{eq:bg:mf:loss} becomes \eqref{eq:bg:fm:loss}. Analytic
    average-velocity matching is thus a strict generalisation of instantaneous-velocity matching, and the fraction of training
    samples drawn with $r \neq \ftime$ is a genuine hyperparameter: at zero it is instantaneous-velocity matching
    \parencite[Tab.~1a]{geng2025meanflow}. The same collapse applies, for a different reason, to the
    consistency-interpolated objective of \autoref{sec:method:alphaflow}, which is how a set of earlier runs
    was found to be mislabelled.
  \item \textbf{The interval is a network input, so the budget is chosen at inference.} As in
    \autoref{sec:method:fm}, $\nfe$ appears only in the sampler. Unlike \autoref{sec:method:fm}, small $\nfe$ is
    no longer paid for in discretisation error, because the field being queried is already the
    exact average over the step. The budget mechanism of RQ1 is therefore a property of what is
    learned.
  \item \textbf{The objective differentiates through the network, which constrains the
    architecture.} Everything inside the closure that \eqref{eq:bg:mf:jvp} differentiates must
    be forward-mode differentiable. This is not a performance detail: it decides how a visual
    encoder may be attached (\autoref{sec:method:backbone}).
\end{itemize}

Two implementation-level choices from \parencite[Sec.~4.3]{geng2025meanflow} are adopted unchanged and are
recorded here because they are not cosmetic. The loss is reweighted adaptively,
$w = (\|\Delta\|_{2}^{2}+c)^{-p}$ applied under a stop-gradient to the squared error $\Delta$,
which normalises the gradient scale across the very different error magnitudes that $r=\ftime$
and $r\ll\ftime$ samples produce; and the time pair is drawn from a logit-normal distribution
rather than uniformly \parencite{esser2024scaling}, with the larger draw assigned to $\ftime$ and the
smaller to $r$. A
consequence of the first is that the reported training loss is pinned near its ceiling by
construction and must never be read as a convergence signal --- the raw squared error is the
quantity to monitor. \srcnote{Both are implemented in
\texttt{flow\_matcher\_v3\_meanflow/models/mf\_diffusion.py} (\texttt{\_adaptive},
\texttt{\_sample\_tau\_pair}); the convergence-signal warning is this work's own, from the
Gen3v6 fix-2/fix-3 logs.}

\subsubsection*{Implementation: the start-anchored identity}

The network takes two times. This work anchors the query at the \emph{start} of the interval,
$(\vect{x}_{r},r,h)$ with $h=\ftime-r$, because that is the point the sampler queries; the
source \parencite[Eq.~(6)]{geng2025meanflow} anchors at the end and integrates in the opposite time
direction. Differentiating $(\ftime-r)\afield = \int_{r}^{\ftime}\vfield$ with respect to $r$ at
fixed $\ftime$, and following the trajectory so that $\dd\vect{x}_{r}/\dd r = \vfield$, gives the
start-anchored form of \eqref{eq:bg:mf:identity},
\begin{equation}\label{eq:method:engine:mfid}
  \afield(\vect{x}_{r},r,\ftime)
  = \vfield(\vect{x}_{r},r) + h\,\frac{\dd}{\dd r}\afield(\vect{x}_{r},r,\ftime),
  \qquad h = \ftime - r,
\end{equation}
which is \eqref{eq:bg:mf:identity} under the change of variables, with the sign flip carried by
$\dd h/\dd r = -1$. The training target and loss are then
\begin{align}
  \afield\sidx{tgt}
    &= \sg\!\left[\;\vect{v} + h\,
       \mathrm{JVP}\!\left(\afield_{\theta};\,
         (\vect{x}_{r},r,h),\,(\vect{v},+1,-1)\right)\right],
       \qquad \vect{v} = \vect{x}_{1}-\vect{x}_{0},
       \label{eq:method:engine:mftgt}\\
  \mathcal{L}(\theta)
    &= \E\Big[\,
        \varrho\big(\|\afield_{\theta}(\vect{x}_{r},r,h) - \afield\sidx{tgt}\|_{2}^{2}\big)
      + \varrho\big(\|\vfield_{\theta}(\vect{x}_{r},r,h) - \sg(\vect{v})\|_{2}^{2}\big)\Big],
       \label{eq:method:engine:mfloss}
\end{align}
where $\varrho(e) = e\,/\,\sg\big((e+c)^{p}\big)$ is the adaptive reweighting of
\autoref{sec:method:meanflow} and the squared errors are summed over the plan before reweighting.
Four properties of \eqref{eq:method:engine:mftgt}--\eqref{eq:method:engine:mfloss} are choices of
this implementation:

\begin{itemize}
  \item \textbf{The $\vect{x}$-tangent is the analytic velocity $\vect{x}_{1}-\vect{x}_{0}$, not a
    predicted one.} Substituting a network-predicted velocity there turns this model into a
    different one, and the resulting ablation would no longer be controlled.
  \item \textbf{$(\ftime,r)$ are two independent draws} on the transport axis, the larger
    assigned to $\ftime$; drawing $r$ conditionally on $\ftime$ starves the large-$h$ queries,
    which is precisely where one- and two-step sampling lives.
  \item \textbf{A fraction of samples is drawn with $r = \ftime$}, i.e.\ $h=0$, at which
    \eqref{eq:method:engine:mftgt} reduces to $\afield\sidx{tgt} = \vect{v}$ and the objective
    becomes \eqref{eq:method:engine:fm}. These anchor the field.
  \item \textbf{The instantaneous head is a full second loss, not a stabiliser}, and it is
    \emph{discarded at sampling time} --- only $\afield_{\theta}$ is integrated. Mixing the two
    heads at sampling time produced step-to-step jitter.
\end{itemize}
\srcnote{\texttt{flow\_matcher\_v3\_meanflow/models/mf\_diffusion.py:353--473}
(\texttt{\_sample\_tau\_pair}, \texttt{\_adaptive}, \texttt{\_p\_losses\_meanflow}); the sign
check against the source's opposite time direction is recorded in that function's docstring.}

Sampling takes $\nfe$ interval jumps of size $h=1/\nfe$, passing both the position and the
interval:
\begin{equation}\label{eq:method:engine:mfsample}
  \vect{x}_{\ftime_{k+1}} = \vect{x}_{\ftime_{k}}
    + \tfrac{1}{\nfe}\,\afield_{\theta}\!\left(\vect{x}_{\ftime_{k}},\;\ftime_{k},\;h{=}\tfrac{1}{\nfe},\;\vect{c}\right),
  \qquad \ftime_{k} = \tfrac{k}{\nfe},
\end{equation}
which at $\nfe = 1$ is exactly the one-step form \eqref{eq:bg:mf:onestep}. The transport time is
never frozen: the network was trained on $\afield_{\theta}(\cdot,\ftime,h)$ and querying it at a
fixed $\ftime$ turns it into a biased $h$-only function, out of the distribution it was fitted
on. Where an adaptive solver is used instead of the Euler scheme, the interval passed at each
sub-stage is the remaining distance to that macro step's declared endpoint, which keeps every
query inside the $\{\ftime,h\ge 0,\ \ftime+h\le 1\}$ domain the model was trained on.
\srcnote{\texttt{mf\_diffusion.py:189--300}; the guardrail and the sub-stage interval rule are
this work's own, from the Gen3v6 U3/U8 audits.}

\subsection{Consistency-Interpolated Average-Velocity Matching}\label{sec:method:alphaflow}

The third objective \parencite{zhang2025alphaflow} also learns the average velocity, but interpolates
its target by a consistency step ratio $\alpha$ between the instantaneous velocity and a stop-gradient
prediction of the network itself. With the architecture matched and the ratio annealed to a final value
$\alpha\sidx{end}>0$, it matches analytic average-velocity matching at low step budgets and
falls behind it at high budgets in both environments where it was measured, and it is reported as an
ablation (\autoref{sec:res:fewstep}, \autoref{sec:disc:negative}).

The target replaces the differentiated term of \eqref{eq:method:engine:mftgt} by the network's own
average velocity, evaluated after travelling a fraction $\alpha$ of the interval at the analytic
velocity \parencite[Def.~1, Alg.~1]{zhang2025alphaflow}. In the start-anchored notation of
\autoref{sec:method:meanflow}, with $\vect{v} = \vect{x}_{1}-\vect{x}_{0}$ and the step
$\delta = \alpha h$,
\begin{equation}\label{eq:method:engine:aftgt}
  \afield\sidx{tgt}
  = \alpha\,\vect{v}
  + (1-\alpha)\,\sg\!\left(\afield_{\theta}\!\left(\vect{x}_{r}+\delta\,\vect{v},\;r+\delta,\;h-\delta\right)\right),
\end{equation}
which is composed into the same reweighted loss \eqref{eq:method:engine:mfloss}. The two endpoints are
exact: at $\alpha = 1$ the target is $\vect{v}$ and the objective is \eqref{eq:method:engine:fm}, and as
$\alpha \to 0$ a first-order expansion of the network term returns the analytic target
\eqref{eq:method:engine:mftgt} \parencite[Thm.~1]{zhang2025alphaflow}. During training $\alpha$ follows
a sigmoid curriculum from one to $\alpha\sidx{end}$
\parencite[Alg.~2]{zhang2025alphaflow}, and a fraction of each batch is drawn with $h = 0$, where
\eqref{eq:method:engine:aftgt} reduces to $\vect{v}$.
\srcnote{\texttt{flow\_matcher\_v3\_alphaflow/models/af\_diffusion.py:637--720}
(\texttt{\_p\_losses\_alphaflow}, $\delta = \alpha h$, the network term under \texttt{no\_grad}) and
\texttt{:440--470} (the sigmoid schedule); upstream
\texttt{PAPERS/auxiliary\_papers/DGM/AlphaFlow\_src/algos/training.tex}.}

Two of its properties matter here. The step ratio appears only in training, and at one
setting the two targets coincide term for term: a checkpoint trained at that setting is analytic
average-velocity matching,
whatever its directory is called, and a set of this work's earlier runs was mislabelled that way. And
its target is a finite difference closed with the network's own stop-gradient output, so it inherits
a fraction of that network's error: the objective contracts error at a rate set by $\alpha$ and
admits self-consistent fields that are wrong. The analytic target of \eqref{eq:method:engine:mftgt}
has no such term. Both this collapse and the $r=\ftime$ collapse of \autoref{sec:method:meanflow} are
extreme hyperparameter values that silently change which model is measured, and both are checked by
assertion.

\subsection{A Common Interface}\label{sec:method:interface}

The inherited system is organised around two calls: one that turns an iterate into a proposed
mean, and one that loops over $\nfe$ steps and applies the constraint mechanism. Every generative model
in this thesis implements exactly those two calls, and the projection code, the candidate plans,
the conditioning, the normalisation and the closed loop are shared verbatim. Written as the next iterate $\vect{\mu}^{k}$ that the projection of
\autoref{sec:method:projection} acts on, the substitution is
\begin{equation}\label{eq:method:engine:iface}
  \vect{\mu}^{k}
  = \begin{cases}
      \dfrac{\beta_{k}\sqrt{\bar{\alpha}_{k-1}}}{1-\bar{\alpha}_{k}}\,\hat{\vect{x}}^{0}
        + \dfrac{(1-\bar{\alpha}_{k-1})\sqrt{\alpha_{k}}}{1-\bar{\alpha}_{k}}\,\vect{x}^{k}
        & \text{diffusion (baseline),\; \eqref{eq:bg:ddpm:mean}}\\[10pt]
      \vect{x}_{\ftime_{k}} + \tfrac{1}{\nfe}\,
        \vfield_{\theta}\!\left(\vect{x}_{\ftime_{k}},\ftime_{k},\vect{c}\right)
        & \text{instantaneous velocity}\\[6pt]
      \vect{x}_{\ftime_{k}} + \tfrac{1}{\nfe}\,
        \afield_{\theta}\!\left(\vect{x}_{\ftime_{k}},\ftime_{k},h{=}\tfrac{1}{\nfe},\vect{c}\right)
        & \text{average velocity,}
    \end{cases}
\end{equation}
with $\ftime_{k} = k/\nfe$ and the reverse index of the first line replaced by the forward index
of the other two. At sampling time, the three lines are the difference between the models.
\srcnote{\texttt{flow\_matcher\_v3/models/diffusion.py:125--158} and
\texttt{flow\_matcher\_v3\_meanflow/models/mf\_diffusion.py:189--300}, against
\texttt{aux\_repo/dpcc/diffuser/models/diffusion.py:119--203}.}

\subsection{Backbone and Conditioning}\label{sec:method:backbone}

One temporal U-Net at one parameter count is used by every generative model, including the
baseline. For the vision-conditioned task, a visual encoder produces a fixed-width latent, and the
backbone never sees an image.

\subsubsection*{The trajectory backbone}

The backbone is a one-dimensional temporal U-Net in the lineage of
\parencite{janner2022planning}: the plan is treated as a length-$H$ signal with $d$ channels,
and convolution runs along \emph{time}, not along the state dimension. Channel widths follow
$\left(c_{0},\dots,c_{L}\right) = \left(d,\;w,\;2w,\;4w,\;8w\right)$ for a base width $w$, with
three stride-two down/up stages and skip concatenations. The horizon is zero-padded up to the next
multiple of eight before the network and cropped back afterwards; $H = 8$ needs no padding.

Every level is a pair of residual blocks. Writing $\vect{h}\in\R^{c\sidx{in}\times H}$ for a block's
input, $\mathcal{B}$ for a $\left(\text{Conv1d}\to\text{GroupNorm}\to\text{Mish}\right)$ unit and
$\vect{z}$ for the block's conditioning vector,
\begin{equation}\label{eq:method:backbone:block}
  \operatorname{Block}(\vect{h},\vect{z})
  = \mathcal{B}_{2}\Big(\mathcal{B}_{1}(\vect{h}) + \big(\matr{W}\sidx{emb}\,\varphi(\vect{z})\big)\vect{1}_{H}\trans\Big)
    + \matr{P}\vect{h},
\end{equation}
with $\varphi$ a pointwise nonlinearity, $\matr{P}$ a $1\times1$ channel-matching convolution, and
$\vect{1}_{H}\trans$ broadcasting the projected embedding along the horizon. The conditioning
enters as a per-channel additive bias that is constant in time.

\subsubsection*{The conditioning vector}

For the state-based environments $\vect{z}$ is the transport-time embedding alone,
$\vect{z} = \phi_{\ftime}(\ftime)$, a sinusoidal positional encoding followed by a two-layer MLP;
conditioning on the current observation is done by inpainting
(\autoref{sec:method:formalisation}).

For the vision-conditioned environment, the visual encoder is that of D3IL \parencite{jia2024towards},
which adopts the encoder of Diffusion Policy \parencite[Sec.~3.2]{chi2023diffusion}. Each of the two
RGB views, at $96\times96$ pixels, has its own ResNet-18 \parencite{he2016deep}, trained end-to-end
with the generative model. Its global average pooling is replaced by spatial-softmax pooling
\parencite{mandlekar2021matters} over $32$ keypoints, followed by a linear layer to $64$ features, and
its batch normalisation by group normalisation \parencite{wu2018group}. The two feature vectors are
concatenated to a per-frame descriptor, and the descriptor is averaged over the
$T\sidx{win}$-frame observation window, giving one latent per plan,
\begin{equation}\label{eq:method:backbone:latent}
  \vect{e}
  = \frac{1}{T\sidx{win}}\sum_{i=1}^{T\sidx{win}}
    \Big[\,E\sidx{view1}(\matr{I}^{(1)}_{i})\;;\;E\sidx{view2}(\matr{I}^{(2)}_{i})\,\Big]
  \;\in\;\R^{128}.
\end{equation}
$\vect{e}$ is the entire interface between perception and generation: the input channels of the
U-Net are the $d$ channels of the plan, and the perception stack is the same for every generative
model.

The latent conditions the backbone by \emph{feature-wise conditional biasing}
\parencite[Sec.~3]{perez2017film}. It is projected by its own MLP and concatenated with the time
embedding,
\begin{equation}\label{eq:method:backbone:concat}
  \vect{z} = \big[\,\phi_{\ftime}(\ftime)\;;\;\phi\sidx{vis}(\vect{e})\,\big] \in \R^{2w},
\end{equation}
so that it reaches every residual block through the additive bias of
\eqref{eq:method:backbone:block}.
\srcnote{\texttt{models/unet1d\_temporal\_cond.py:53--82, 205--235};
\texttt{mix\_visual\_aligning/models/visual\_unet.py:64--97}. Name mapping for the configuration flag
and run tags: Auxiliary/Naming/TRANSLATION\_20260914\_dev\_jargon\_to\_scientific.md \S5.}

\subsubsection*{Two-time models}

The two average-velocity objectives need $\afield_{\theta}(\cdot,\ftime,h)$, so the backbone takes a
second time input. It is embedded by its own MLP and added to the transport-time embedding, before
the visual latent is concatenated:
\begin{equation}\label{eq:method:backbone:twotime}
  \vect{z} = \big[\,\phi_{\ftime}(\ftime) + \phi_{h}(h)\;;\;\phi\sidx{vis}(\vect{e})\,\big].
\end{equation}
This is the only structural change the two-time models make to the backbone.

The analytic target \eqref{eq:method:engine:mftgt} differentiates the network with a forward-mode
\ac{JVP} in $(\vect{x}_{r},r,h)$. The latent $\vect{e}$ depends on none of these inputs, so the encoder
is evaluated once, outside the differentiated function, and $\vect{e}$ enters it as a constant with
zero tangent. The two ResNet towers are therefore not differentiated by the \ac{JVP}, and they are
still trained end-to-end through the loss.
\srcnote{Methodology\_Sources/AUX\_visual\_aligning\_env.md \S3.1;
\texttt{mix\_visual\_aligning/models/visual\_unet\_twotime.py:1--45}.}

% =============================================================================
%  CONSTRAINT PROJECTION
% =============================================================================
\section{Constraint Projection}\label{sec:method:projection}

Both projection methods solve the same program on the same feasible set with the same solver. Iterate
projection applies it to the sampler's current iterate; endpoint projection applies it to the plan the
sampler is predicted to reach. Plans generated with neither are unguided.

\subsection{The Projection Problem}\label{sec:method:proj}

\subsubsection*{The feasible set and the projection}

Let $\mathcal{Z} \subset \R^{n}$ collect the plans that satisfy the test-time state and action
constraints, and let $\mathcal{Z}_{f} \subseteq \mathcal{Z}$ additionally impose the dynamics
model. Following \parencite[Eq.~(9)]{romer2025diffusion}, the \emph{model-based} projection of a plan is the
program
\begin{align}
  \proj{\mathcal{Z}_{f}}(\vect{x})
    &= \arg\min_{\tilde{\vect{x}}\in\mathcal{Z}} \;\|\vect{x}-\tilde{\vect{x}}\|_{2}^{2}
    \label{eq:bg:mpc:proj} \\
  \text{s.t.}\quad & \vect{s}_{t'+1|t} = f\!\left(\vect{s}_{t'|t},\vect{a}_{t'|t}\right),
    \qquad \forall t' \in \{t,\dots,t+H\},
    \label{eq:bg:mpc:dyn}
\end{align}
with projection cost $c_{\mathcal{Z}_{f}}(\vect{x}) = \|\vect{x}-\proj{\mathcal{Z}_{f}}(\vect{x})\|_{2}^{2}$.
The projection is \emph{goal-independent}: it knows nothing about the task, so it can make a
plan admissible but cannot make it more successful. Both projection methods therefore
change the learned plan as little as possible. The projection is
also a \emph{nonlinear} program whenever the geometric
constraints are non-convex, as circular keep-out regions make them; the solver used throughout
this thesis is \ac{SLSQP}
\parencite{kraft1988software,virtanen2020scipy}, inherited from \parencite{romer2025diffusion}.

\subsubsection*{The program}

With the plan flattened to $\vect{x}\in\R^{n}$, \eqref{eq:bg:mpc:proj} is instantiated as a
quadratic objective with linear equality rows for the dynamics, linear inequality rows for the
bounds, and one quadratic inequality per horizon step per obstacle:
\begin{align}
  \proj{\mathcal{Z}_{f}}(\vect{x})
    = \arg\min_{\vect{z}\in\R^{n}}\;\; &
       \tfrac{1}{2}\,\vect{z}\trans\matr{Q}\vect{z} - \vect{x}\trans\matr{Q}\vect{z}
       \;=\; \tfrac{1}{2}\left\|\vect{z}-\vect{x}\right\|_{\matr{Q}}^{2} + \text{const}
       \label{eq:method:proj:obj}\\
  \text{s.t.}\quad &
       \matr{A}\vect{z} = \vect{b}
       &&\text{(dynamics rows, \eqref{eq:bg:mpc:dyn})}
       \label{eq:method:proj:eq}\\
    & \matr{C}\vect{z} \le \vect{d}
       &&\text{(state / action bounds, halfspaces)}
       \label{eq:method:proj:ineq}\\
    & \vect{z}_{i}\trans\matr{P}_{j}\vect{z}_{i} + \vect{\ell}_{j}\trans\vect{z}_{i}
       \le \nu_{j}
       &&\text{(obstacle $j$ at horizon step $i$),}
       \label{eq:method:proj:obs}
\end{align}
with $\matr{Q}$ diagonal --- the identity unless a subset of dimensions is weighted --- and
$\vect{z}_{i}$ the $i$-th horizon block of $\vect{z}$; $(\matr{P}_{j},\vect{\ell}_{j},\nu_{j})$
are the quadratic, linear and constant coefficients of obstacle $j$. The rows \eqref{eq:method:proj:obs} are
what makes $\mathcal{Z}_{f}$ non-convex, so \eqref{eq:method:proj:obj}--\eqref{eq:method:proj:obs}
is a \ac{NLP} and is solved by SLSQP \parencite{kraft1988software,virtanen2020scipy}, warm
started at $\vect{x}$ itself, with box bounds inherited from the normalisation. The first
dynamics block is overwritten with the measured current state, so the projection cannot move the
point the plan is anchored at. The returned projection cost is
$c_{\mathcal{Z}_{f}}(\vect{x}) = \tfrac{1}{2}\|\proj{\mathcal{Z}_{f}}(\vect{x})-\vect{x}\|_{\matr{Q}}^{2}$,
which is the quantity the cumulative-cost selection rule of \autoref{sec:method:selection} sums.
Each candidate plan is a separate program.
\srcnote{\texttt{aux\_repo/dpcc/diffuser/sampling/projection.py:70--155} and its Gen-local copies;
the constraint families are \texttt{SafetyConstraints}, \texttt{DynamicConstraints},
\texttt{ObstacleConstraints} in the same file.}

\subsubsection*{The dynamics rows, written out}

The Euler model \eqref{eq:method:dpcc:euler} enters \eqref{eq:method:proj:eq} as one equality row per
integrated channel per horizon step. Writing $z_{i,d}$ for dimension $d$ of horizon block $i$, a
channel pair $(d\sidx{pos},d\sidx{act})$ contributes
\begin{equation}\label{eq:method:proj:deriv}
  z_{i,d\sidx{pos}} + t\sidx{s}\,z_{i,d\sidx{act}} - z_{i+1,d\sidx{pos}} = 0,
  \qquad i = 0,\dots,H-1,
\end{equation}
plus one row pinning $z_{0,d\sidx{pos}}$ to the measured current state --- the only row whose
right-hand side changes at every control step. Two details of \eqref{eq:method:proj:deriv} are easy
to get wrong and are therefore stated. First, the projector operates in \emph{normalised}
coordinates, so each row is rescaled by the per-dimension data range: with
$\Delta_{d} = z_{\textrm{max},d} - z_{\textrm{min},d}$ and
$\Sigma_{d} = z_{\textrm{max},d} + z_{\textrm{min},d}$, the implemented row is
\begin{equation}\label{eq:method:proj:derivnorm}
  \Delta_{d\sidx{pos}}\,z_{i,d\sidx{pos}}
  + t\sidx{s}\,\Delta_{d\sidx{act}}\,z_{i,d\sidx{act}}
  - \Delta_{d\sidx{pos}}\,z_{i+1,d\sidx{pos}}
  = -\,t\sidx{s}\,\Sigma_{d\sidx{act}} ,
\end{equation}
which is \eqref{eq:method:proj:deriv} conjugated by the normaliser; forgetting the affine offset
$\Sigma$ silently biases every dynamics row. Second, the row count is the number of
\emph{channel pairs}, and that number is a modelling decision, not a constant: it is four for the
two-dimensional manipulator and six for both three-dimensional environments, because each of them
anchors a commanded and a measured position in three axes
(\autoref{tab:embodiments}).
\srcnote{\texttt{aux\_repo/dpcc/diffuser/sampling/projection.py:344--402}
(\texttt{DynamicConstraints.build\_matrices}); the pair lists are in
\texttt{mix\_visual\_aligning\_test/eval\_mix\_visual\_aligning.py:272--279} and
\texttt{mix\_uav\_test/eval\_mix\_uav.py:1265--1266}.}

\subsubsection*{The activation schedule}

The projection is not applied at every step. A threshold $\eta\in[0,1]$ selects the portion of
the transport near the data end at which it runs. In the inherited descending-index sampler that
is the condition $k \le \eta\nfe$; in this work's ascending-index transport samplers it is
\begin{equation}\label{eq:method:proj:gate}
  k \;\ge\; \left\lfloor (1-\eta)\,\nfe \right\rfloor
  \qquad\text{or}\qquad k = \nfe-1 ,
\end{equation}
the second clause forcing the terminal step so that a plan is never returned unprojected. The
two forms agree up to the rounding of $\eta\nfe$, and the number of activated steps is
\begin{equation}\label{eq:method:proj:nact}
  n\sidx{act}(\nfe,\eta) = \max\!\left\{\nfe - \left\lfloor(1-\eta)\nfe\right\rfloor,\;1\right\}.
\end{equation}
$\eta = 0$ therefore does not mean ``no projection'': it means the single terminal projection,
i.e.\ project-after-sampling. This is the post-processing baseline of \parencite[Sec.~6.1]{romer2025diffusion},
and writing it as a special case of \eqref{eq:method:proj:gate} rather than as a separate code
path is what keeps the two comparable.
\srcnote{\texttt{flow\_matcher\_v3/models/diffusion.py:174--183} carries the derivation of the
guarded form and the record of what it changes relative to the earlier unguarded expression.}

\subsection{Per-Step Projection (DPCC)}\label{sec:method:dpcc}

Per-step projection is the constraint mechanism of \ac{DPCC} \parencite{romer2025diffusion}.

\subsubsection*{The constrained denoising step}

The inherited algorithm is \eqref{eq:bg:ddpm:reverse} with a projection inserted between its
steps. Its justification is a guidance argument: introducing a binary feasibility variable
$O$ with likelihood
$p(O{=}1\mid\vect{x},k) \propto \exp\!\big(-\tfrac{1}{2\sigma_{k}^{2}}d(\vect{x},\mathcal{Z}_{f})^{2}\big)$,
where $d$ is the Euclidean distance to the feasible set, a first-order expansion of
$\log p(O\mid\vect{x},k)$ about the learned mean --- the same expansion that underlies classifier
guidance \parencite{dhariwal2021diffusion} --- gives a guidance term that, for closed convex
$\mathcal{Z}_{f}$, is exactly the displacement to the projection. The conditioned reverse step
is therefore
\begin{equation}\label{eq:method:dpcc:thm}
  p_{\theta}\!\left(\vect{x}^{k-1}\mid\vect{x}^{k},k,O{=}1\right)
  \approx \Normal\!\left(\proj{\mathcal{Z}_{f}}\!\left(\vect{\mu}_{\theta}^{k}\right),\;
    \sigma_{k}^{2}\Id{n}\right)
\end{equation}
\parencite[Theorem~1]{romer2025diffusion}. Because \eqref{eq:method:dpcc:thm} projects the
\emph{mean} and then adds noise, the drawn sample can still leave $\mathcal{Z}_{f}$; the
algorithm that is actually run therefore reorders the two operations, projecting \emph{after}
the noise is added:
\begin{equation}\label{eq:method:dpcc:step}
  \vect{x}^{k-1}
  = \proj{\mathcal{Z}_{f}}\!\left(
      \vect{\mu}_{\theta}\!\left(\vect{x}^{k},k,\vect{c}\right) + \sigma_{k}\vect{\epsilon}^{k}
    \right),
  \qquad \vect{\epsilon}^{k}\sim\Normal(\vect{0},\Id{n}),
\end{equation}
which gives deterministic constraint satisfaction at the steps where it is applied
\parencite[Eq.~16]{romer2025diffusion}. \eqref{eq:method:dpcc:step} is the single equation this
thesis replaces: \autoref{sec:method:engine} substitutes the generative model that produces
$\vect{\mu}_{\theta}$, and \autoref{sec:method:hardflow} substitutes the operator applied to it.

\subsubsection*{What is projected}

This method projects the \emph{current iterate}. Under \eqref{eq:method:proj:gate} with a large $\eta$
that iterate is still far from the data end of the transport, so
\eqref{eq:method:proj:obj}--\eqref{eq:method:proj:obs} is posed on an object that is not yet a
plan: the geometric rows \eqref{eq:method:proj:obs} are being asked about positions that the
sampler has not finished producing. This is the single structural difference from endpoint
projection, and \autoref{sec:disc:interpretation} builds on it.

\subsubsection*{Constraint tightening}

The dynamics $f$ used inside \eqref{eq:bg:mpc:dyn} is an approximation, with mismatch bounded
by $\|\vect{w}_{t}\|_{2}\le\gamma$. Feasibility of the \emph{predicted} plan therefore does not
imply feasibility of the executed one. The inherited remedy tightens the predicted-state
constraints by the mismatch ball,
\begin{equation}\label{eq:method:dpcc:tighten}
  \tilde{\mathcal{S}}_{t+1} = \mathcal{S}_{t+1}\ominus\mathcal{B}_{\gamma},
\end{equation}
with $\ominus$ the Minkowski difference and $\mathcal{B}_{\gamma}$ the $\ell_{2}$ ball of radius
$\gamma$; sampling against $\tilde{\mathcal{Z}}_{f}$ then guarantees that all executed states
satisfy the original constraints \parencite[Theorem~2]{romer2025diffusion}. The bound $\gamma$ is
estimated empirically from unconstrained rollouts. \emph{Tightened geometry} in this thesis
means \eqref{eq:method:dpcc:tighten} applied with a larger $\gamma$ than the baseline's; it is varied
as an experimental factor.

\subsection{Endpoint Projection}\label{sec:method:hardflow}

Endpoint projection \parencite{li2025hardflow} is re-derived here on this work's own constraint set.

\subsubsection*{Sampling as optimal control}

A control $\vect{u}_{\ftime}$ is added to the learned velocity field, and the constraints are imposed on
the final sample only \parencite[Problem~1]{li2025hardflow}:
\begin{equation}\label{eq:method:hf:oc}
  \min_{\vect{u}_{[0,1]}}\; C(\vect{x}_{1})
    + \lambda\sidx{oc}\int_{0}^{1}\tfrac{1}{2}\left\|\vect{u}_{\ftime}\right\|_{2}^{2}\,\dd\ftime
  \quad\text{s.t.}\quad
  \dot{\vect{x}}_{\ftime} = \vfield_{\theta}(\vect{x}_{\ftime},\ftime,\vect{c}) + \vect{u}_{\ftime},
  \qquad \vect{x}_{1}\in\mathcal{Z}_{f},
\end{equation}
starting from the drawn noise $\vect{x}_{0}$. $C$ is an optional terminal cost, and the effort term keeps
the controlled sampler close to the learned one.

\subsubsection*{Predicting the final sample}

For an affine interpolation path the expected final sample given an intermediate state follows in
closed form from the state and the marginal velocity \parencite[Lemma~1]{li2025hardflow}. For the linear
path \eqref{eq:bg:fm:path} the expectations of the final sample and of the initial noise are
\begin{equation}\label{eq:method:hf:post}
  \mathcal{M}_{\ftime}(\vect{x}) = \E\!\left[\vect{x}_{1}\mid\vect{x}_{\ftime}=\vect{x}\right]
    = \vect{x} + (1-\ftime)\,\vfield(\vect{x},\ftime),
  \qquad
  \mathcal{W}_{\ftime}(\vect{x}) = \E\!\left[\vect{x}_{0}\mid\vect{x}_{\ftime}=\vect{x}\right]
    = \vect{x} - \ftime\,\vfield(\vect{x},\ftime),
\end{equation}
with $\vect{x} = \ftime\,\mathcal{M}_{\ftime}(\vect{x}) + (1-\ftime)\,\mathcal{W}_{\ftime}(\vect{x})$.
At $\ftime=1$ the prediction is the state itself.

\subsubsection*{One problem per sampling step}

The method discretises \eqref{eq:method:hf:oc} on the Euler grid of \eqref{eq:bg:fm:euler}, splits it into
one problem per step in which the final sample is replaced by its prediction
\eqref{eq:method:hf:post}, and changes the decision variable from the next state to the predicted final
sample, so that the constraint set is not distorted by the network \parencite[Problems~2--5]{li2025hardflow}.
Inverting the prediction with one fixed-point step gives the algorithm it runs
\parencite[Algorithm~1]{li2025hardflow}. With $\Delta\ftime = 1/\nfe$, a guided step $k$ reads
\begin{align}
  \bar{\vect{x}}_{k+1} &= \vect{x}_{\ftime_{k}}
    + \Delta\ftime\,\vfield_{\theta}(\vect{x}_{\ftime_{k}},\ftime_{k},\vect{c}),
    \label{eq:method:hf:ref}\\
  \bar{\vect{x}}_{1} &= \bar{\vect{x}}_{k+1}
    + (1-\ftime_{k+1})\,\vfield_{\theta}(\bar{\vect{x}}_{k+1},\ftime_{k+1},\vect{c}),
    \label{eq:method:hf:pred}\\
  \hat{\vect{x}}_{1}^{\star} &= \arg\min_{\hat{\vect{x}}_{1}\in\mathcal{Z}_{f}}\;
    C(\hat{\vect{x}}_{1})
    + \frac{\lambda\sidx{oc}\,\ftime_{k+1}^{2}}{2\,\Delta\ftime}\left\|\hat{\vect{x}}_{1}-\bar{\vect{x}}_{1}\right\|_{2}^{2},
    \label{eq:method:hf:solve}\\
  \vect{x}_{\ftime_{k+1}} &= \bar{\vect{x}}_{k+1} + \ftime_{k+1}\left(\hat{\vect{x}}_{1}^{\star}-\bar{\vect{x}}_{1}\right).
    \label{eq:method:hf:pull}
\end{align}
\eqref{eq:method:hf:ref} is an unguided Euler step, \eqref{eq:method:hf:pred} predicts its final sample,
\eqref{eq:method:hf:solve} corrects the prediction, and \eqref{eq:method:hf:pull} ---
$\ftime_{k+1}\hat{\vect{x}}_{1}^{\star} + (1-\ftime_{k+1})\,\mathcal{W}_{\ftime_{k+1}}(\bar{\vect{x}}_{k+1})$
written out --- pulls the correction back into the sampler, damped by $\ftime_{k+1}$. At the final step
$\ftime_{\nfe} = 1$, so the output equals $\hat{\vect{x}}_{1}^{\star}\in\mathcal{Z}_{f}$ and satisfies the
constraints \parencite[Prop.~1]{li2025hardflow}. The optimisation runs only in the later sampling
steps, where the prediction \eqref{eq:method:hf:post} is accurate \parencite{li2025hardflow}.
\eqref{eq:method:hf:ref}--\eqref{eq:method:hf:pull} are the steps of
\parencite[Alg.~1]{li2025hardflow}, written in the notation of this thesis; the source shows the
sequence of corrected predictions on the obstacle-avoidance task used here
\parencite[Fig.~11]{li2025hardflow}.

\subsubsection*{The configuration used here}

The program \eqref{eq:method:hf:solve} is built from the constraint list of iterate projection
(\autoref{sec:method:proj}), and the terminal cost is zero, $C\equiv 0$. The objective is then a
quadratic proximity term whose minimiser does not depend on its weight, so
\begin{equation}\label{eq:method:hf:ours}
  \hat{\vect{x}}_{1}^{\star} = \proj{\mathcal{Z}_{f}}\!\left(\bar{\vect{x}}_{1}\right),
\end{equation}
the projection of \autoref{sec:method:proj} applied to the predicted final sample. Constraints and
solver are those of iterate projection, and only the point of projection differs; the factor
$\ftime_{k+1}$ in \eqref{eq:method:hf:pull} damps the early steps. Each guided step evaluates the
network twice, in \eqref{eq:method:hf:ref} and \eqref{eq:method:hf:pred}.
\autoref{tab:hardflow-setup} lists the configuration.

\begin{table}[htpb]
  \caption[Configuration of endpoint projection]{Configuration of endpoint projection in this thesis.}
  \label{tab:hardflow-setup}
  \centering
  \begin{tabular}{ll}
    \toprule
    plan length               & $8$ steps \\
    actions executed per plan & $1$ \\
    step budget $\nfe$        & $1$--$20$ \\
    solver                    & SLSQP \parencite{kraft1988software} \\
    dynamics rows             & Euler, $t\sidx{s}=1$, \eqref{eq:method:dpcc:euler} \\
    terminal cost $C$         & none \\
    activation                & second half of the steps, $\eta = 0.5$ \\
    \bottomrule
  \end{tabular}
\end{table}
\srcnote{\texttt{flow\_matcher\_v3\_meanflow/sampling/hardflow\_projection.py:1--40} (mode
\texttt{hardflow\_new}, steps \eqref{eq:method:hf:ref}--\eqref{eq:method:hf:pull} with $C\equiv0$);
\texttt{config/avoiding-d3il.py:318}; per-timestep replanning in
\texttt{aux\_repo/dpcc/scripts/eval.py:203--241}.}

\subsection{Steps at Which Projection Guides Sampling}\label{sec:method:degenerate}

The activation convention \eqref{eq:method:proj:gate} of endpoint projection can leave the terminal step
as the only active one. At the terminal time the lookahead vanishes, the pull-back replaces the sample
by the projection of its prediction, and no later step exists for the network to react to. What runs is
then projection applied to a finished sample.

\subsubsection*{The guiding step count}

Of the $n\sidx{act}$ steps counted by \eqref{eq:method:proj:nact}, one is the terminal step, which does
not guide the sampler. The \emph{guiding step count} is
\begin{equation}\label{eq:method:degen:ngen}
  n\sidx{guide}(\nfe,\eta)
  \;=\; n\sidx{act}(\nfe,\eta) - 1
  \;=\; \max\!\left\{\nfe - \left\lfloor(1-\eta)\nfe\right\rfloor,\;1\right\} - 1 .
\end{equation}
It depends on the step budget and the threshold together, so two configurations with the same step
budget and different thresholds can differ in it. Its value decides what a result can show:

\begin{center}
\begin{tabularx}{\linewidth}{lLL}
  \toprule
  $n\sidx{guide}$ & what runs & what a result supports \\
  \midrule
  $0$     & projection of the final sample only        & no statement about endpoint projection \\
  $1$     & one guided step before the final one       & no attributable effect \\
  $\ge 2$ & guided steps, each followed by a later step & a statement about the method \\
  \bottomrule
\end{tabularx}
\end{center}

The \emph{first lookahead}, the distance $1-\ftime_{k_{0}+1}$ remaining at the first guiding step with
$k_{0} = \nfe - n\sidx{act}$, is reported alongside it. It measures how far the prediction of the final
sample extrapolates at that step.
\srcnote{\texttt{flow\_matcher\_v3\_meanflow/sampling/hardflow\_projection.py:584--646}
(\texttt{hardflow\_step\_budget}, \texttt{hardflow\_regime}); artefact token \texttt{n\_genuine}.}

\subsubsection*{Reporting rule}

Results with $n\sidx{guide} = 0$ are excluded from every statement about endpoint projection; results
with $n\sidx{guide} = 1$ are marked as carrying no attributable effect; neither is averaged with other
results. The endpoint-projection results at small step budgets and $\eta = 0.5$ with
$n\sidx{guide} \le 1$ are listed in \autoref{sec:res:constraints:degenerate}.

\subsection{Candidate Selection}\label{sec:method:selection}

\subsubsection*{Candidate plans}

Rather than one plan, $B$ plans are drawn per timestep and one is executed. The inherited
selection rules are \emph{temporal consistency} --- take the candidate closest to the previous
timestep's executed plan on their overlapping segment --- and \emph{cumulative projection
cost} --- take the candidate that \eqref{eq:bg:mpc:proj} altered least, summed over the steps
at which it ran \parencite[Sec.~5.4]{romer2025diffusion}. The third rule in this work's implementation is random selection,
which is the convention the two rules were proposed against.

\subsubsection*{The three selection rules}

Written out, the three rules are
\begin{equation}\label{eq:method:sel:rules}
  i\sidx{rand}(t) \sim \Unif\{1,\dots,B\},
  \quad
  i\sidx{tc}(t) = \arg\min_{j}\left\|\vect{x}^{j}_{t:t+H-1|t} - \vect{x}^{i(t-1)}_{t:t+H-1|t-1}\right\|_{2},
  \quad
  i\sidx{cc}(t) = \arg\min_{j}\sum_{k} c_{\mathcal{Z}_{f}}\!\left(\tilde{\vect{x}}^{k,j}\right),
\end{equation}
with $\tilde{\vect{x}}$ the pre-projection iterate. With a single candidate, $B=1$, all three rules
return the same index and are the same computation, so \autoref{ch:results} compares them only for
$B>1$.

% =============================================================================
%  ENVIRONMENTS
% =============================================================================
\section{Environments}\label{sec:method:envs}

Obstacle avoidance is the
benchmark of \ac{DPCC}, taken over with its setup unchanged, and is where the new generative models
and projection methods are first established against the baseline. Vision-conditioned alignment and the
quadrotor scenes were built for this thesis and test them further, in a new observation modality and a
new embodiment. Where several environments appear together, they are named D3IL-avoiding,
D3IL-aligning, and, per scene, UAV-corridor, UAV-pillars and UAV-s-curve.

The three environments share the plan layout of \autoref{tab:embodiments} and differ in their
observations, their dynamics rows and their constraint sets. Each constraint below enters the program
\eqref{eq:method:proj:obj}--\eqref{eq:method:proj:obs} as one family of rows, acts on the planned
position $\vect{p}_{i}$ of every horizon step $i$, and is shrunk by the tightening
\eqref{eq:method:dpcc:tighten} in the tightened geometries. Scene dimensions are configuration values and
are listed with the experimental setup in \autoref{sec:setup:tasks}.

Two constraint forms recur. A halfspace, given by a boundary line and the side to be excluded, is the
linear row
\begin{equation}\label{eq:method:env:halfspace}
  \vect{n}\trans\vect{p}_{i} \le d ,
\end{equation}
with unit normal $\vect{n}$ and offset $d$. A circular keep-out region with centre $\vect{o}$ and radius
$\rho$ is the quadratic row
\begin{equation}\label{eq:method:env:disk}
  \left\|\vect{p}_{i}-\vect{o}\right\|_{2}^{2} \ge \rho^{2},
\end{equation}
which makes the feasible set non-convex.
\srcnote{Halfspace rows: \texttt{utils/constraints\_helpers.py} (\texttt{formulate\_halfspace\_constraints});
quadratic rows: \texttt{aux\_repo/dpcc/diffuser/sampling/projection.py:410--462}
(\texttt{ObstacleConstraints}).}

\subsection{Obstacle Avoidance}\label{sec:method:env:avoiding}

The obstacle-avoidance task of D3IL \parencite{jia2024towards} is the benchmark of \ac{DPCC}
\parencite[Sec.~6.1]{romer2025diffusion}. It is simulated in MuJoCo
\parencite{todorov2012mujoco}: a seven-joint Franka Emika Panda carries a cylindrical rod as its end
effector and moves in a plane \parencite[Sec.~3.3]{jia2024towards}. The state stacks the commanded and the measured planar
end-effector position, $\vect{s} = (\vect{p}\sidx{des},\vect{p})\in\R^{4}$, the action is a planar
increment $\Delta\vect{p}\in\R^{2}$, and the four dynamics rows are \eqref{eq:method:dpcc:euler}. The
test-time constraints are halfspaces \eqref{eq:method:env:halfspace} and circular regions
\eqref{eq:method:env:disk} on the measured position, none of which the demonstrations respect, together
with a box on the action given by the range of the demonstrated actions
\parencite[Sec.~6.1]{romer2025diffusion}. The baseline runs in the configuration of \ac{DPCC} --- the
same plan horizon, number of diffusion steps, action weight and number of candidate plans --- on the
same constraint geometries.
\srcnote{Setup parity checked 2026-09-14: \texttt{aux\_repo/dpcc/config/avoiding-d3il.py:22--27, 69}
(horizon 8, \texttt{n\_diffusion\_steps} 20, \texttt{action\_weight} 10, plan \texttt{batch\_size} 4)
against \texttt{config/avoiding-d3il.py:318--323} and \texttt{\_mpc\_batch} (default 4); the two
\texttt{projection\_eval.yaml} differ only in \texttt{write\_to\_file} and \texttt{seeds}, so the
constraint geometries are identical. Candidate count varied as a factor elsewhere is v3's protocol.}

\subsection{Vision-Conditioned Alignment}\label{sec:method:env:aligning}

The alignment task of D3IL \parencite{jia2024towards} asks the robot to push a box into a target
pose. It runs in MuJoCo; its task files, demonstrations, camera placement and vision pipeline are
D3IL's; the constrained-control environment
around them was built for this thesis --- the plan and its dynamics rows, the constraint set, and the
coupling of the visual encoder to the generative network.
\srcnote{Methodology\_Sources/AUX\_visual\_aligning\_env.md \S1 (``the task is D3IL's, the vision
pipeline is D3IL's, and the pairing is ours'').}

It uses the controller stack of the manipulator; what changes is the plan and the observation. The action is a three-dimensional increment $\Delta\vect{p}\in\R^{3}$, the state
half of the observation is $(\vect{p}\sidx{des},\vect{p})\in\R^{6}$ --- commanded and measured
end-effector position --- and the plan width is therefore $d = 9$ per horizon step, laid out as
$(\vect{a},\vect{p}\sidx{des},\vect{p})$. The six dynamics rows of
\eqref{eq:method:proj:deriv} are the pairs
\begin{equation}\label{eq:method:dep:varows}
  \left(d\sidx{pos},d\sidx{act}\right) \in
  \{(3,0),(4,1),(5,2)\}\;\cup\;\{(6,0),(7,1),(8,2)\},
\end{equation}
i.e.\ the commanded triple and the measured triple both integrate the same action, which is
\eqref{eq:method:dpcc:euler} written in three dimensions.
\srcnote{\texttt{mix\_visual\_aligning\_test/eval\_mix\_visual\_aligning.py:272--279}.}

Perception sits above the plan: two RGB views are encoded to a $128$-dimensional latent that
conditions the backbone through the feature-wise bias of \eqref{eq:method:backbone:concat}, and the
backbone never sees an image (\autoref{sec:method:backbone}). Below the setpoint, the control chain is
that of the state-based benchmark.
\srcnote{Methodology\_Sources/AUX\_visual\_aligning\_env.md \S2.}

Its constraint set is a workspace box on the planned end-effector position,
$\vect{p}\sidx{min}\le\vect{p}_{i}\le\vect{p}\sidx{max}$, and a box on the action given by the range of
the demonstrated actions; the table carries no walls, so the task has no halfspaces.
\srcnote{\texttt{config/visual\_aligning\_eval.yaml}: \texttt{workspace\_bounds} per geometry,
\texttt{action\_bounds: 'auto'}.}

\subsection{Quadrotor Scenes}\label{sec:method:env:uav}

The vehicle is the Skydio X2 from MuJoCo Menagerie \parencite{zakka2022menagerie}, copied verbatim and
patched only to spawn level and to drop an unused sensor block and keyframe, and simulated in MuJoCo
\parencite{todorov2012mujoco}. Everything around the vehicle was built for this thesis: the scenes,
their constraint sets, the expert demonstrations and the control loop of \autoref{sec:method:deployment}.
Four static scenes include the vehicle: an empty scene, a corridor between two parallel walls, an
s-curve of two offset corridor segments, and a field of six cylindrical pillars. Walls and pillars are passive rigid bodies and do not change the dynamics of the
vehicle. The plan stacks action, commanded position, measured position and measured velocity
(\autoref{tab:embodiments}), and its six dynamics rows anchor the commanded and the measured position as
in \eqref{eq:method:dep:varows}.

The constraint set of each scene combines four families:
\begin{itemize}
  \item a workspace box $\vect{p}\sidx{min}\le\vect{p}_{i}\le\vect{p}\sidx{max}$ bounding extent and
    altitude;
  \item wall halfspaces \eqref{eq:method:env:halfspace} on the corridor and on the s-curve, where each
    wall applies only over the $x$-interval $[x\sidx{a},x\sidx{b}]$ of its segment,
    \begin{equation}\label{eq:method:env:switched}
      p_{i,x}\in[x\sidx{a},x\sidx{b}] \;\Rightarrow\; \vect{n}\trans\vect{p}_{i} \le d ;
    \end{equation}
  \item keep-out disks \eqref{eq:method:env:disk} for the pillars, for the free ends of the corridor
    walls and for the inner corners of the s-curve;
  \item a box on the action given by the range of the demonstrated actions.
\end{itemize}
Walls and disks are inflated by the vehicle radius $r\sidx{b}$, so a plan for the vehicle's centre keeps
its body clear. The empty scene carries no constraints and serves as the unconstrained reference.
\srcnote{\texttt{config/uav\_projection.yaml}: \texttt{corridor\_hg}, \texttt{pillars\_hg},
\texttt{s\_curve\_hg} (\texttt{planning\_inflation.r\_drone}, per-segment \texttt{x\_active}); model
provenance in \texttt{Methodology\_Sources/AUX\_uav\_model\_and\_scenes.md}.}

\subsubsection*{Expert demonstrations}

The demonstrations are generated in two layers: a geometric reference trajectory, and its execution
by the cascaded geometric tracking controller of \autoref{sec:method:geometric} in the MuJoCo scene.

The reference of an episode is built from a chain of waypoints. The waypoints fix on which side of each
obstacle the path passes; in scenes with several such classes, the classes are cycled across episodes.
Consecutive waypoints are joined by straight segments, and every non-collinear interior waypoint is
replaced by a circular arc of radius at most $r\sidx{f}$ tangent to both adjacent segments, with the
radius reduced where the arc would otherwise use more than half of either segment. The result is a
path $\vect{p}\sidx{path}:[0,L\sidx{path}]\to\R^{3}$ parametrised by arc length, travelled with one cosine
speed profile over the episode duration $T\sidx{ep}$,
\begin{equation}\label{eq:method:env:uavref}
  \vect{p}\sidx{des}(t) = \vect{p}\sidx{path}\big(s(t)\big),
  \qquad
  s(t) = \frac{L\sidx{path}}{2}\left(1-\cos\frac{\pi t}{T\sidx{ep}}\right),
  \qquad 0\le t\le T\sidx{ep},
\end{equation}
so the reference is at rest only at its start and its end. Its velocity and acceleration,
\begin{equation}\label{eq:method:env:uavrefderiv}
  \dot{\vect{p}}\sidx{des} = \vect{p}\sidx{path}'(s)\,\dot{s},
  \qquad
  \ddot{\vect{p}}\sidx{des} = \vect{p}\sidx{path}'(s)\,\ddot{s} + \vect{p}\sidx{path}''(s)\,\dot{s}^{2},
\end{equation}
are the feed-forward terms of \eqref{eq:method:dep:outer}; $\vect{p}\sidx{path}''$ vanishes on the segments
and is the centripetal term on the arcs. The altitude and the duration are drawn at random per
episode, together with the start and end points in the empty scene, a small lateral offset of the start
in the centre lane of the corridor, and a small lateral offset of the whole path in the s-curve.

The vehicle starts level and at rest, and the controller tracks
\eqref{eq:method:env:uavref}--\eqref{eq:method:env:uavrefderiv} at the physics step $\Delta t\sidx{phys}$. An episode is discarded if the vehicle is in
contact with an obstacle for more than a scene-specific fraction of its steps, or if it falls below a
minimum altitude. The accepted rollouts are recorded at the planning interval, every $n\sidx{dec}$-th
physics step of \eqref{eq:method:dep:multirate}, as
\begin{equation}\label{eq:method:env:uavdata}
  \vect{o}_{j} = \left(\vect{p}\sidx{des}(t_{j}),\;\vect{p}(t_{j}),\;\dot{\vect{p}}(t_{j})\right),
  \qquad
  \vect{a}_{j} = \vect{p}\sidx{des}(t_{j+1}) - \vect{p}\sidx{des}(t_{j}),
  \qquad
  t_{j} = j\,n\sidx{dec}\,\Delta t\sidx{phys},
\end{equation}
the form of \eqref{eq:method:dep:setpoint}--\eqref{eq:method:dep:obs}. The action is the increment of
the reference, not an output of the controller: the generative model learns to plan the reference,
and at deployment the same controller tracks it.
\srcnote{Reference: \texttt{uav\_expert\_data\_collect/trajectories.py} (\texttt{pillar\_path},
\texttt{corridor\_path}, \texttt{s\_curve\_scene\_path}, \texttt{empty\_path}),
\texttt{uav\_env\_test/trajectories.py:67--90} (\texttt{traverse\_line}) and \texttt{:114--}
(\texttt{blended\_path}, circular arcs, centripetal term). Execution and rejection:
\texttt{uav\_expert\_data\_collect/generator.py} (\texttt{run\_trial}, \texttt{\_build\_traj\_and\_init},
\texttt{SCENE\_MAX\_CONTACT\_FRACTION}, \texttt{Z\_FLOOR\_MARGIN}); homotopy cycling in
\texttt{collect.py}. Recording: \texttt{dataset\_writer.py} (\texttt{\_downsample}, stride 3;
\texttt{actions = np.diff(targets)}). The per-episode offset on \texttt{targets} is not stated: it is
added before the difference and cancels in it, and \texttt{obs} is built from the un-offset
reference. Values of $r\sidx{f}$, thresholds and ranges belong to \autoref{sec:setup:data}.}

% =============================================================================
% CHAPTERS 5-8: BONE ONLY -- DELIBERATELY NOT WRITTEN
% Nothing here is settled. Do not fill these in until the run queue in
% Data_Analysis/DA_Result_Curated_MD/NOTEBOOK_20260829_key_headlines.md clears.
% =============================================================================

